\documentclass[letterpaper,11pt]{article}

\usepackage[utf8]{inputenc}
\usepackage[T1]{fontenc}
\usepackage[brazilian]{babel}
\usepackage{graphicx}
%\usepackage{fullpage}
\usepackage{pdfpages}
\usepackage{color}
\usepackage[colorlinks=true,urlcolor=blue,citecolor=black]{hyperref}
\usepackage{url}
\usepackage[font=footnotesize,labelfont=bf]{caption}
%nome completo para o apêndice
\usepackage[title]{appendix}
\usepackage{float}
%\usepackage{parskip}
%para código
\usepackage{listings}
%para matemática
\usepackage{amsmath}

\linespread{1.1}
\setlength{\parindent}{30pt}
\setlength{\emergencystretch}{3em}


%abertura
\title{\LARGE A Lightning Network do Bitcoin:\\
	\Large Pagamentos Instantâneos Off-Chain e Escaláveis}
\author{
		Joseph Poon\\
		\small\href{mailto:joseph@lightning.network}
			{\nolinkurl{joseph@lightning.network}}
	\and
		Thaddeus Dryja\\
		\small\href{mailto:rx@awsomnet.org}
			{\nolinkurl{rx@awsomnet.org}}
	}
\date{\today\\\small RASCUNHO Versão 0.5.9.2 \\ \small Tradução para o Português (Brasil)}
\begin{document}

\maketitle

\begin{abstract}

O protocolo Bitcoin pode abranger todo o volume global de transações financeiras
de todos os sistemas de pagamento eletrônico atuais, sem que um único terceiro
custodiante mantenha os fundos ou que se exija dos participantes mais do que um
computador com uma conexão de banda larga. É proposto um sistema descentralizado
no qual as transações são enviadas por uma rede de canais de micropagamento
(também conhecidos como canais de pagamento ou canais de transação) cuja
transferência de valor ocorre fora da blockchain. Se as transações Bitcoin
puderem ser assinadas com um novo tipo de sighash que resolva a maleabilidade,
essas transferências podem ocorrer entre partes que não confiam umas nas outras
ao longo da rota de transferência, por meio de contratos que, em caso de
participantes não cooperativos ou hostis, são exigíveis via transmissão na
blockchain do Bitcoin, através de uma série de timelocks decrescentes.

\end{abstract}

\section{O Problema de Escalabilidade da Blockchain Bitcoin}

A blockchain do Bitcoin\cite{nakamoto} é muito promissora como livro-razão
distribuído, mas a blockchain como plataforma de pagamento, por si só, não pode
abranger todo o comércio mundial em um futuro próximo. A blockchain é um
protocolo de fofoca (gossip protocol), no qual todas as modificações de estado
do livro-razão são transmitidas a todos os participantes. É por meio desse
``protocolo de fofoca'' que se chega a um consenso sobre o estado, isto é, sobre
os saldos de todos. Se cada nó da rede Bitcoin precisar saber de cada uma das
transações que ocorrem globalmente, isso pode prejudicar substancialmente a
capacidade da rede de abranger todas as transações financeiras globais. Em vez
disso, seria desejável abranger todas as transações de uma forma que não
sacrifique a descentralização e a segurança que a rede oferece.

A rede de pagamentos Visa atingiu o pico de 47.000 transações por segundo (tps)
em sua rede durante o feriado de 2013\cite{visa}, e atualmente processa em média
centenas de milhões de transações por dia. Atualmente, o Bitcoin suporta menos
de 7 transações por segundo, com um limite de bloco de 1 megabyte. Se usarmos
uma média de 300 bytes por transação Bitcoin e supusermos blocos de tamanho
ilimitado, uma capacidade equivalente ao pico de transações da Visa de 47.000
tps exigiria blocos de quase 8 gigabytes a cada dez minutos, em média.
Continuamente, isso resultaria em mais de 400 terabytes de dados por ano.

Claramente, atingir hoje uma capacidade similar à da Visa na rede Bitcoin não é
viável. Nenhum computador doméstico do mundo consegue operar com esse nível de
largura de banda e armazenamento. Se o Bitcoin viesse a substituir todos os
pagamentos eletrônicos no futuro, e não apenas a Visa, o resultado seria o
colapso total da rede Bitcoin ou, na melhor das hipóteses, uma centralização
extrema dos nós e mineradores nas poucas entidades que pudessem arcar com isso.
Essa centralização anularia aspectos da descentralização da rede que tornam o
Bitcoin seguro, já que é justamente a capacidade de entidades validarem a cadeia
que permite ao Bitcoin garantir a exatidão e a segurança do livro-razão.

Ter menos validadores devido a blocos maiores implica não apenas em menos
indivíduos garantindo a exatidão do livro-razão, mas também em menos entidades
capazes de validar a blockchain como parte do processo de mineração, o que
estimula a centralização da mineração. Blocos extremamente grandes, como no caso
acima de 8 gigabytes a cada 10 minutos em média, implicariam que apenas algumas
poucas entidades seriam capazes de validar blocos. Isso cria uma grande
possibilidade de que entidades acabem confiando em partes centralizadas. Ter
partes privilegiadas e confiáveis cria uma armadilha social, na qual a parte
central não agirá em favor do indivíduo (problema do agente-principal),
incluindo, por exemplo, rentismo via cobrança de taxas mais altas para mitigar o
incentivo a agir desonestamente. Em casos extremos, isso se manifesta como
indivíduos enviando fundos a custodiantes centralizados de confiança que têm a
custódia total dos fundos dos clientes. Tais arranjos, como ocorrem hoje, criam
sério risco de contraparte. Um pré-requisito para impedir que esse tipo de
centralização ocorra exige a capacidade de o Bitcoin ser validado por um único
computador de consumidor em uma conexão doméstica de banda larga. Ao garantir
que a validação completa possa ocorrer de forma barata, os nós e mineradores
Bitcoin serão capazes de evitar centralização extrema e confiança, o que assegura
taxas de transação extremamente baixas.

Embora seja possível que a Lei de Moore continue indefinidamente, e que a
capacidade computacional para que os nós computem economicamente blocos de
múltiplos gigabytes venha a existir no futuro, isso não é uma certeza.

Para atingir muito mais do que 47.000 transações por segundo usando o Bitcoin,
é necessário conduzir as transações fora da própria blockchain do Bitcoin. Seria
ainda melhor se a rede Bitcoin suportasse um número quase ilimitado de
transações por segundo, com taxas extremamente baixas para micropagamentos.
Muitos micropagamentos podem ser enviados sequencialmente entre duas partes para
viabilizar pagamentos de qualquer tamanho. Os micropagamentos permitiriam o
desempacotamento, menor necessidade de confiança e a comoditização de serviços,
como pagamentos por megabyte para serviço de internet. Para poder atender a
esses casos de uso de micropagamento, no entanto, seria necessário reduzir
drasticamente a quantidade de transações que acabam sendo transmitidas para a
blockchain global do Bitcoin.

Embora seja possível escalar em um nível pequeno, é absolutamente impossível
lidar com uma grande quantidade de micropagamentos na rede ou abranger todas as
transações globais. Para que o Bitcoin tenha sucesso, é preciso ter confiança
de que, se ele se tornar extremamente popular, suas vantagens atuais decorrentes
da descentralização continuarão existindo. Para que as pessoas hoje acreditem
que o Bitcoin funcionará amanhã, o Bitcoin precisa resolver a questão dos
efeitos de centralização do tamanho do bloco; blocos grandes criam,
implicitamente, custodiantes de confiança e taxas significativamente mais
altas.

\section{Uma Rede de Canais de Micropagamento Pode Resolver a Escalabilidade}

\begin{quote}
	``Se uma árvore cai na floresta e não há ninguém por perto para ouvir,
	ela faz algum som?''
\end{quote}

A citação acima questiona a relevância de eventos não observados \textemdash se
ninguém ouve a árvore cair, se ela fez som ou não é irrelevante. De forma
similar, na blockchain, se apenas dois participantes se importam com uma
transação cotidiana recorrente, não é necessário que todos os demais nós da rede
Bitcoin saibam dessa transação. É preferível, em vez disso, ter apenas o mínimo
de informação na blockchain. Ao adiar o ato de contar ao mundo inteiro sobre
cada transação, fazendo a liquidação líquida da relação em uma data posterior,
os usuários do Bitcoin podem realizar muitas transações sem inchar a blockchain
ou criar confiança em uma contraparte centralizada. Uma estrutura efetivamente
sem confiança pode ser obtida usando timelocks como componente do consenso
global.

Atualmente, a solução para micropagamentos e escalabilidade é transferir as
transações para um custodiante, no qual se confia em terceiros para guardar as
moedas e atualizar saldos com outras partes. Confiar em terceiros para guardar
todos os fundos cria risco de contraparte e custos de transação.

Em vez disso, usando uma rede desses canais de micropagamento, o Bitcoin pode
escalar para bilhões de transações por dia com o poder computacional disponível
em um computador desktop moderno atual. Enviar muitos pagamentos dentro de um
mesmo canal de micropagamento permite enviar grandes quantias de fundos para
outra parte de forma descentralizada. Esses canais não são uma rede
intermediária confiável construída por cima do Bitcoin. Eles são transações
Bitcoin reais.

Os canais de micropagamento\cite{wikicontracts}\cite{bitcoinjmicropay} criam uma
relação entre duas partes para atualizar saldos perpetuamente, adiando o que é
transmitido à blockchain para uma única transação que liquida o saldo total
entre essas duas partes. Isso permite que as relações financeiras entre duas
partes sejam adiadas, sem confiança, para uma data posterior, sem risco de
inadimplência da contraparte. Os canais de micropagamento usam transações
Bitcoin reais, optando apenas por adiar a transmissão à blockchain de modo que
ambas as partes possam garantir seu saldo atual na blockchain; não se trata de
uma rede sobreposta confiável \textemdash os pagamentos em canais de
micropagamento são bitcoins reais comunicados e trocados fora da cadeia.

\subsection{Canais de Micropagamento Não Requerem Confiança}

Assim como a antiga questão sobre a árvore que cai na floresta e se faz som, se
todas as partes concordam que a árvore caiu às 14:45 da tarde, então a árvore
realmente caiu às 14:45 da tarde. De forma similar, se ambas as contrapartes
concordam que o saldo atual dentro de um canal é 0,07 BTC para Alice e 0,03 BTC
para Bob, então esse é o saldo verdadeiro. Entretanto, sem criptografia, surge
um problema interessante: se uma contraparte discorda do saldo atual dos fundos
(ou do horário em que a árvore caiu), é a palavra de um contra a palavra do
outro. Sem assinaturas criptográficas, a blockchain não saberá quem é dono do
quê.

Se o saldo no canal for 0,05 BTC para Alice e 0,05 BTC para Bob, e o saldo
após uma transação for 0,07 BTC para Alice e 0,03 BTC para Bob, a rede precisa
saber qual conjunto de saldos está correto. As transações da blockchain
resolvem esse problema usando o livro-razão da blockchain como um sistema de
carimbo de tempo. Ao mesmo tempo, é desejável criar um sistema que não use
ativamente esse sistema de carimbo de tempo, a menos que absolutamente
necessário, já que ele pode se tornar custoso para a rede.

Em vez disso, ambas as partes podem se comprometer a assinar uma transação e
não transmiti-la. Assim, se Alice e Bob colocarem fundos em um endereço
multiassinatura 2-de-2 (em que é necessário o consentimento de ambas as partes
para criar gastos), eles podem concordar sobre o estado atual do saldo. Alice
e Bob podem concordar em criar um reembolso a partir dessa transação 2-de-2
para si mesmos, 0,05 BTC para cada um. Esse reembolso \textit{não} é
transmitido na blockchain. Qualquer uma das partes pode transmiti-lo, mas
podem optar por manter essa transação consigo, sabendo que podem resgatar os
fundos quando se sentirem confortáveis. Ao adiar a transmissão dessa
transação, podem optar por alterar esse saldo em uma data futura.

Para atualizar o saldo, ambas as partes criam um novo gasto a partir do
endereço multiassinatura 2-de-2, por exemplo, 0,07 para Alice e 0,03 para Bob.
Sem o projeto adequado, porém, surge o problema de carimbo de tempo: não se
sabe qual gasto é correto: o novo gasto ou o reembolso original.

A restrição quanto ao carimbo de tempo e datas, contudo, não é tão complexa
quanto a ordenação total de todas as transações como ocorre na blockchain do
Bitcoin. No caso dos canais de micropagamento, apenas dois estados são
necessários: o saldo atual correto e quaisquer saldos antigos depreciados.
Existe apenas um único saldo atual correto, e possivelmente muitos saldos
antigos depreciados.

Portanto, é possível, no Bitcoin, conceber um script de modo que todas as
transações antigas sejam invalidadas e apenas a nova transação seja válida.
A invalidação é imposta por um script de saída Bitcoin e por transações
dependentes que forçam a outra parte a entregar todos os seus fundos à
contraparte do canal. Ao confiscar todos os fundos como penalidade para serem
entregues à outra parte, todas as transações antigas são, assim, invalidadas.

Esse processo de invalidação pode existir através de um processo de consenso
do canal em que, se ambas as partes concordam com os estados atuais do
livro-razão (e na construção de novos estados), então o saldo real é
atualizado. O saldo é refletido na blockchain apenas quando uma única parte
discorda. Conceitualmente, esse sistema não é uma rede sobreposta independente;
é mais um adiamento de estado sobre o sistema atual, já que a aplicação ainda
ocorre na própria blockchain (embora adiada para datas e transações futuras).

\subsection{Uma Rede de Canais}

Assim, os canais de micropagamento criam apenas uma relação entre duas partes.
Exigir que todos criem canais com todos os outros não resolve o problema de
escalabilidade. A escalabilidade do Bitcoin pode ser obtida usando uma grande
rede de canais de micropagamento.

Se presumirmos uma grande rede de canais sobre a blockchain do Bitcoin, e que
todos os usuários do Bitcoin estão participando desse grafo tendo ao menos um
canal aberto na blockchain do Bitcoin, é possível criar uma quantidade quase
infinita de transações dentro dessa rede. As únicas transações transmitidas à
blockchain do Bitcoin prematuramente são aquelas com contrapartes de canal não
cooperativas.

%TODO: definir hashlock e timelock
Ao encumbrar (onerar) as saídas das transações Bitcoin com um hashlock e
timelock, a contraparte do canal será incapaz de simplesmente roubar fundos e
os bitcoins podem ser trocados sem risco de roubo direto pela contraparte.
Além disso, ao usar timeouts escalonados, é possível enviar fundos por meio de
vários intermediários em uma rede sem o risco de roubo dos fundos por esses
intermediários.

\section{Canais de Pagamento Bidirecionais}

Os canais de micropagamento permitem o simples adiamento da transmissão de um
estado de transação para um momento posterior. Os contratos são impostos
criando uma responsabilidade de uma parte de transmitir transações antes ou
depois de certas datas. Se a blockchain é um sistema descentralizado de
carimbo de tempo, é possível usar relógios como componente do consenso
descentralizado\cite{lamportpaxos} para determinar a validade dos dados, assim
como apresentar estados como método para ordenar
eventos\cite{lamportclocks}.

Ao criar janelas de tempo dentro das quais certos estados podem ser
transmitidos e posteriormente invalidados, é possível criar contratos
complexos usando scripts de transação Bitcoin. Há trabalhos anteriores sobre
Canais de Micropagamento em estrela
(Hub-and-Spoke)\cite{akselrod}\cite{akselrod2}\cite{todd} (e redes de canais de
pagamento confiáveis\cite{amikopay}\cite{impulse}) buscando construir uma rede
em estrela atualmente. Entretanto, o canal de micropagamento bidirecional da
Lightning Network requer o soft-fork de maleabilidade descrito no Apêndice A
para permitir escalabilidade quase infinita, mitigando os riscos de
inadimplência dos nós intermediários.

Ao encadear múltiplos canais de micropagamento, é possível criar uma rede de
caminhos de transações. Os caminhos podem ser roteados usando um sistema
similar ao BGP, e o remetente pode designar um caminho específico para o
destinatário. Os scripts de saída são onerados por um hash, gerado pelo
destinatário. Ao divulgar a entrada (preimage) desse hash, a contraparte do
destinatário poderá puxar os fundos ao longo da rota.

\subsection{O Problema de Atribuição de Culpa na Criação do Canal}

Para participar dessa rede de pagamentos, é necessário criar um canal de
micropagamento com outro participante da rede.

\subsubsection{Criando uma Transação de Financiamento Não Assinada}

Uma Transação de Financiamento inicial do canal é criada, na qual uma ou ambas
as contrapartes do canal financiam as entradas dessa transação. Ambas as
partes criam as entradas e saídas dessa transação, mas não assinam a
transação.

A saída dessa Transação de Financiamento é um único script multiassinatura
2-de-2 com ambos os participantes do canal, daqui em diante chamados de Alice
e Bob. Ambos os participantes não trocam assinaturas para a Transação de
Financiamento até que tenham criado gastos a partir dessa saída 2-de-2,
reembolsando o valor original aos respectivos financiadores. O propósito de
não assinar a transação permite gastar a partir de uma transação que ainda
não existe. Se Alice e Bob trocassem as assinaturas da Transação de
Financiamento sem antes poderem transmitir gastos a partir dela, os fundos
poderiam ficar travados para sempre caso Alice e Bob não cooperassem (ou
outras perdas de moedas poderiam ocorrer através de cenários de refém, em que
se paga pela cooperação da contraparte).

Alice e Bob trocam as entradas para financiar a Transação de Financiamento
(para saber quais entradas serão usadas para determinar o valor total do
canal) e trocam uma chave a ser usada para assinar mais tarde. Essa chave é
usada para a saída 2-de-2 da Transação de Financiamento; ambas as assinaturas
são necessárias para gastar a partir da Transação de Financiamento, ou seja,
tanto Alice quanto Bob precisam concordar em gastar a partir dela.

\subsubsection{Gastando a partir de uma Transação Não Assinada}

A Lightning Network usa uma transação SIGHASH\_NOINPUT para gastar a partir
dessa saída 2-de-2 da Transação de Financiamento, pois é necessário gastar a
partir de uma transação cujas assinaturas ainda não foram trocadas.
SIGHASH\_NOINPUT, implementado via soft-fork, garante que as transações possam
ser gastas antes de serem assinadas por todas as partes, já que, sem novos
flags de sighash, seria necessário assinar para obter um ID de transação.
Sem o SIGHASH\_NOINPUT, as transações Bitcoin não podem ser gastas antes de
poderem ser transmitidas \textemdash é como se não fosse possível elaborar um
contrato sem antes pagar a outra parte. O SIGHASH\_NOINPUT resolve esse
problema. Veja o Apêndice A para mais informações e implementação.

Sem o SIGHASH\_NOINPUT, não é possível gerar um gasto a partir de uma
transação sem trocar assinaturas, já que gastar a Transação de Financiamento
requer o ID da transação como parte da assinatura na entrada do filho. Um
componente do ID da Transação é a assinatura do pai (Transação de
Financiamento), de modo que ambas as partes precisam trocar suas assinaturas
da transação pai antes que o filho possa ser gasto. Como uma ou ambas as
partes precisam conhecer as assinaturas do pai para gastar a partir dele,
isso significa que uma ou ambas as partes seriam capazes de transmitir o pai
(Transação de Financiamento) antes que o filho sequer exista. O
SIGHASH\_NOINPUT contorna isso permitindo que o filho seja gasto sem assinar
a entrada. Com o SIGHASH\_NOINPUT, a ordem das operações é:
\begin{enumerate}
	\item Criar o pai (Transação de Financiamento)
	\item Criar os filhos (Transações de Compromisso e todos os gastos a
		partir das transações de compromisso)
	\item Assinar os filhos
	\item Trocar as assinaturas dos filhos
	\item Assinar o pai
	\item Trocar as assinaturas do pai
	\item Transmitir o pai na blockchain
\end{enumerate}

Não é possível transmitir o pai (Etapa 7) até que a Etapa 6 esteja completa.
Ambas as partes não deram sua assinatura para gastar a partir da Transação
de Financiamento até a etapa 6. Além disso, se uma parte falhar durante a
Etapa 6, o pai pode ser gasto (tornando-se a transação pai) ou as entradas
da transação pai podem ser gastas em duplicidade (de modo que todo esse
caminho de transação seja invalidado).

\subsubsection{Transações de Compromisso: Construção Não Exequível}

Após a criação da Transação de Financiamento não assinada (e não transmitida),
ambas as partes assinam e trocam uma Transação de Compromisso inicial. Essas
Transações de Compromisso gastam a partir da saída 2-de-2 da Transação de
Financiamento (pai). Entretanto, apenas a Transação de Financiamento é
transmitida na blockchain.

Como a Transação de Financiamento já entrou na blockchain, e a saída é uma
transação multiassinatura 2-de-2 que requer o acordo de ambas as partes para
ser gasta, as Transações de Compromisso são usadas para expressar o saldo
presente. Se apenas uma Transação de Compromisso 2-de-2 assinada for trocada
entre as partes, ambas terão certeza de que poderão recuperar seu dinheiro
depois que a Transação de Financiamento entrar na blockchain. As partes não
transmitem as Transações de Compromisso na blockchain até que desejem fechar o
saldo atual do canal. Elas fazem isso transmitindo a Transação de Compromisso
atual.

As Transações de Compromisso pagam os saldos atuais respectivos a cada parte.
Uma implementação ingênua (defeituosa) construiria uma transação não
transmitida em que há um gasto 2-de-2 a partir de uma única transação com
duas saídas que devolvem todos os saldos atuais às duas contrapartes do
canal. Isso devolverá todos os fundos à parte original ao criar uma Transação
de Compromisso inicial.

%Diagrama - Transação de financiamento defeituosa
\begin{figure}[H]
	\makebox[\linewidth]{
		\includegraphics[width=1\linewidth]{figures/funding-broken1.pdf}
	}
	\caption{Este diagrama descreve uma transação de financiamento
		ingênua e defeituosa. A Transação de Financiamento (F),
		designada em verde, é transmitida na blockchain após todas as
		demais transações serem assinadas. Todas as outras transações
		que gastam a partir da transação de financiamento ainda não
		são transmitidas, caso as contrapartes desejem atualizar seu
		saldo. Apenas a Transação de Financiamento é transmitida na
		blockchain neste momento.
	}
\end{figure}

Por exemplo, se Alice e Bob concordam em criar uma Transação de Financiamento
com uma única saída 2-de-2 no valor de 1,0 BTC (com contribuição de 0,5 BTC
de cada um), eles criam uma Transação de Compromisso com duas saídas de
0,5 BTC para Alice e Bob. As Transações de Compromisso são assinadas primeiro
e as chaves são trocadas para que qualquer um possa transmitir a Transação de
Compromisso a qualquer momento, condicionado à entrada da Transação de
Financiamento na blockchain. Nesse ponto, as assinaturas da Transação de
Financiamento podem ser trocadas com segurança, já que qualquer das partes é
capaz de resgatar seus fundos transmitindo a Transação de Compromisso.

Esta construção, contudo, falha quando se deseja atualizar o saldo atual.
Para atualizar o saldo, é necessário atualizar os valores de saída da
Transação de Compromisso (a Transação de Financiamento já entrou na
blockchain e não pode ser alterada).

Quando ambas as partes concordam com uma nova Transação de Compromisso e
trocam assinaturas para a nova Transação de Compromisso, qualquer uma das
Transações de Compromisso pode ser transmitida. Como a saída da Transação
de Financiamento só pode ser resgatada uma vez, apenas uma dessas transações
será válida. Por exemplo, se Alice e Bob concordam que o saldo do canal é
agora 0,4 para Alice e 0,6 para Bob, e uma nova Transação de Compromisso é
criada para refletir isso, qualquer uma das Transações de Compromisso pode
ser transmitida. Na prática, não há como restringir qual Transação de
Compromisso é transmitida, já que ambas as partes assinaram e trocaram as
assinaturas para que qualquer saldo possa ser transmitido.

%Diagrama - Transação de financiamento defeituosa
\begin{figure}[H]
	\makebox[\linewidth]{
		\includegraphics[width=1\linewidth]{figures/funding-broken2.pdf}
	}
	\caption{Qualquer uma das Transações de Compromisso pode ser
		transmitida a qualquer momento por qualquer parte; apenas uma
		conseguirá efetivamente gastar a única Transação de
		Financiamento. Isso não pode funcionar porque uma das partes
		não vai querer transmitir a transação mais recente.
	}
\end{figure}

Como qualquer parte pode transmitir a Transação de Compromisso a qualquer
momento, o resultado seria que, após gerada a nova Transação de Compromisso,
aquela que recebe menos fundos tem incentivo significativo para transmitir a
transação que tem valores maiores para si nas saídas da Transação de
Compromisso. Como resultado, o canal seria imediatamente fechado e os fundos
roubados. Portanto, não é possível criar canais de pagamento sob esse modelo.

\subsubsection{Transações de Compromisso: Atribuindo Culpa}

Como qualquer Transação de Compromisso assinada pode ser transmitida na
blockchain, e apenas uma pode ser transmitida com sucesso, é necessário
impedir que Transações de Compromisso antigas sejam transmitidas. Não é
possível revogar dezenas de milhares de transações no Bitcoin, então é
necessário um método alternativo. Em vez de revogação ativa imposta pela
blockchain, é preciso construir o próprio canal de forma similar a uma
Caução de Fidelidade (Fidelity Bond), em que ambas as partes fazem
compromissos, e violações desses compromissos são impostas via penalidades.
Se uma parte violar seu acordo, ela perderá todo o dinheiro do canal.

Para este canal de pagamento, os termos do contrato são que ambas as partes
se comprometem a transmitir apenas a transação mais recente. Qualquer
transmissão de transações mais antigas causará uma violação do contrato, e
todos os fundos são entregues à outra parte como penalidade.

Isso só pode ser imposto se for possível atribuir culpa pela transmissão de
uma transação antiga. Para isso, deve-se ser capaz de identificar de forma
única quem transmitiu uma transação antiga. Isso pode ser feito se cada
contraparte tiver uma Transação de Compromisso identificável de forma única.
Ambas as partes precisam assinar as entradas da Transação de Compromisso pela
qual a outra parte é responsável pela transmissão. Como cada parte tem uma
versão da Transação de Compromisso assinada pela outra, cada parte só pode
transmitir sua própria versão da Transação de Compromisso.

Para a Lightning Network, todos os gastos a partir da saída da Transação de
Financiamento, as Transações de Compromisso, têm duas transações
meio-assinadas. Uma Transação de Compromisso na qual Alice assina e entrega
a Bob (C1b), e outra na qual Bob assina e entrega a Alice (C1a). Essas duas
Transações de Compromisso gastam a partir da mesma saída (Transação de
Financiamento) e têm conteúdos diferentes; apenas uma pode ser transmitida
na blockchain, pois ambos os pares de Transações de Compromisso gastam a
partir da mesma Transação de Financiamento. Qualquer parte pode transmitir
sua Transação de Compromisso recebida assinando sua versão e incluindo a
assinatura da contraparte. Por exemplo, Bob pode transmitir a Transação de
Compromisso C1b, já que recebeu a assinatura de C1b de Alice \textemdash ele
inclui a assinatura de Alice e assina C1b ele mesmo. A transação será um
gasto válido a partir da saída 2-de-2 da Transação de Financiamento, que
exige tanto a assinatura de Alice quanto a de Bob.

%Diagrama - Transação de financiamento defeituosa
\begin{figure}[H]
	\makebox[\linewidth]{
		\includegraphics[width=1\linewidth]{figures/funding-broken3.pdf}
	}
	\caption{Caixas roxas são transações não transmitidas que apenas Alice
		pode transmitir. Caixas azuis são transações não transmitidas
		que apenas Bob pode transmitir. Alice só pode transmitir a
		Transação de Compromisso 1a, Bob só pode transmitir a
		Transação de Compromisso 1b. Apenas uma Transação de
		Compromisso pode ser gasta a partir da saída da Transação de
		Financiamento. A culpa é atribuída, mas qualquer uma ainda
		pode ser gasta sem penalidade.
	}
\end{figure}

Entretanto, mesmo com essa construção, apenas se atribuiu culpa. Ainda não é
possível impor esse contrato na blockchain do Bitcoin. Bob ainda confia em
Alice para não transmitir uma Transação de Compromisso antiga. Neste momento,
ele apenas é capaz de provar que Alice fez isso por meio de uma prova de
transação meio-assinada.

\subsection{Criando um Canal com Revogação de Contrato}

Para poder efetivamente impor os termos do contrato, é necessário construir
uma Transação de Compromisso (junto com seus gastos) na qual seja possível
revogar uma transação. Essa revogação é alcançável usando dados sobre quando
uma transação entra na blockchain e usando a maturidade da transação para
determinar os caminhos de validação.

\subsection{Maturidade de Número de Sequência}

Mark Friedenbach propôs que Números de Sequência possam ser exigíveis via uma
maturidade relativa de blocos da transação pai através de um
soft-fork\cite{seqnum}. Isso permitiria alguma habilidade básica de garantir
alguma forma de bloqueio de tempo baseado em confirmação relativa de blocos
no script de gasto. Além disso, um opcode adicional,
OP\_CHECKSEQUENCEVERIFY\cite{csv} (também conhecido como
OP\_RELATIVECHECKLOCKTIMEVERIFY)\cite{relCLTV}, permitiria habilidades
adicionais, incluindo permitir uma solução paliativa antes de uma solução
mais permanente para resolver a maleabilidade de transações. Uma versão
futura deste artigo incluirá soluções propostas.

Em resumo, o Bitcoin foi lançado com um número de sequência que era exigido
apenas no mempool de transações não confirmadas. O comportamento original
permitia a substituição de transações ao substituir transações no mempool por
novas transações com número de sequência maior. Devido às regras de
substituição de transações, isso não é imposto devido a riscos de ataque de
negação de serviço. Aparentemente, a intenção do número de sequência seria
substituir transações não transmitidas. Entretanto, esse comportamento de
substituição por número de sequência maior é inexequível. Não é possível
garantir que versões antigas de transações tenham sido substituídas no
mempool e que um bloco contenha a versão mais recente da transação. Uma forma
de impor versões de transação fora da cadeia é via compromissos de tempo.

Uma Transação Revogável gasta a partir de uma saída única em que a transação
tem um tipo único de script de saída. A saída desse pai tem dois caminhos de
resgate: o primeiro pode ser resgatado imediatamente, e o segundo só pode ser
resgatado se o filho tiver um número mínimo de confirmações entre as
transações. Isso é alcançado fazendo com que o número de sequência da
transação filho exija um número mínimo de confirmações a partir do pai. Na
essência, esse novo comportamento de número de sequência só permitirá que um
gasto dessa saída seja válido se o número de blocos entre a saída e a
transação que a resgata estiver acima de uma altura de bloco especificada.

Uma transação pode ser revogada com esse comportamento de número de sequência
criando uma restrição com algum número de blocos definido no número de
sequência, o que resultará em o gasto só ser válido depois que o pai tenha
entrado na blockchain por um número definido de blocos. Isso cria uma
estrutura na qual a transação pai com essa saída se torna um depósito
caucionado, atestando que não há revogação. Existe um período de tempo no
qual qualquer um na blockchain pode contestar essa atestação transmitindo um
gasto imediatamente após a transação ser transmitida.

Se alguém deseja permitir transações revogáveis com um atraso de 1000
confirmações, a construção da transação de saída permaneceria como um
multisig 2-de-2:

\begin{lstlisting}
2 <Alice1> <Bob1> 2 OP_CHECKMULTISIG
\end{lstlisting}

Entretanto, a transação filho que gasta conteria um valor de nSequence de
1000. Como essa transação requer a assinatura de ambas as contrapartes para
ser válida, ambas as partes incluem o número de nSequence 1000 como parte
da assinatura. Ambas as partes podem, a seu critério, concordar em criar
outra transação que substitua essa transação sem nenhum número de
nSequence.

Essa construção, um Contrato de Maturidade de Sequência Revogável (Revocable
Sequence Maturity Contract, RSMC), cria dois caminhos, com termos contratuais
muito específicos.

Os termos do contrato são:
\begin{enumerate}
	\item Todas as partes pagam um contrato com uma saída que impõe esse
		contrato
	\item Ambas as partes podem concordar em enviar fundos para algum
		contrato, com algum período de espera (1000 confirmações em
		nosso script de exemplo). Este é o saldo de saída revogável.
	\item Uma ou ambas as partes podem optar por não transmitir (impor) os
		pagamentos até uma data futura; qualquer parte pode resgatar
		os fundos após o período de espera a qualquer momento.
	\item Se nenhuma das partes transmitiu essa transação (resgatou os
		fundos), elas podem revogar os pagamentos acima se, e somente
		se, ambas concordarem em fazê-lo, colocando um novo termo de
		pagamento em uma transação substitutiva. O novo pagamento da
		transação pode ser resgatado imediatamente depois que o
		contrato é divulgado ao mundo (transmitido na blockchain).
	\item Caso o contrato seja divulgado e a nova estrutura de pagamento
		não seja resgatada, os termos de pagamento revogados
		anteriormente podem ser resgatados por qualquer parte (de modo
		que é responsabilidade de qualquer das partes impor os novos
		termos).
\end{enumerate}

A transação filho pré-assinada pode ser resgatada depois que a transação pai
tiver entrado na blockchain com 1000 confirmações, em razão do número
nSequence do filho na entrada que gasta o pai.

Para revogar essa transação filho assinada, ambas as partes simplesmente
concordam em criar outra transação filho com o campo padrão de nSequence
igual a MAX\_INT, que tem comportamento especial permitindo gasto a qualquer
momento.

Esse novo gasto assinado substitui o gasto revogável, desde que o novo gasto
assinado entre na blockchain dentro de 1000 confirmações da entrada do pai
na blockchain. Na prática, se Alice e Bob concordam em monitorar a
blockchain em busca de transmissão incorreta de Transações de Compromisso,
no momento em que a transação é transmitida, eles podem gastar usando a
transação substitutiva imediatamente. Para transmitir o gasto revogável
(transação depreciada), que gasta a partir da mesma saída que a transação
substitutiva, eles devem esperar 1000 confirmações. Desde que ambas as
partes monitorem a blockchain, o gasto revogável nunca entrará na transação
se qualquer das partes preferir a transação substitutiva.

Usando essa construção, qualquer um poderia criar uma transação, não
transmiti-la, e mais tarde criar incentivos para nunca mais transmitir essa
transação no futuro via penalidades. Isso permite que participantes da rede
Bitcoin adiem muitas transações de nunca atingirem a blockchain.

\subsubsection{Timestop}

Para mitigar uma inundação de transações por um atacante malicioso, é
necessária uma ameaça crível de que o ataque falhará.

Greg Maxwell propôs usar um timestop para mitigar uma inundação maliciosa na
blockchain:

\begin{quote}
	Existem muitas maneiras de lidar com isso [risco de inundação] que
	ainda não foram adequadamente exploradas \textemdash por exemplo, o
	relógio pode parar quando os blocos estão cheios, transformando o
	risco de segurança em mais atraso de retenção em caso de um ataque de
	negação de serviço.\cite{gregtimestop}
\end{quote}

Isso pode ser mitigado permitindo que o minerador especifique se o mempool
atual (taxa paga) está sendo inundado com transações. Eles podem entrar com
um valor ``1'' no último bit do número de versão do cabeçalho do bloco. Se o
último bit no cabeçalho do bloco contiver ``1'', então esse bloco não contará
para a maturidade de altura relativa do valor de nSequence, e o bloco é
designado como um bloco congestionado. Existe uma altura de bloco não
congestionada (que é sempre menor que a altura de bloco normal). Essa altura
de bloco é usada para o valor de nSequence, que só conta a maturidade do
bloco (confirmações).

Um minerador pode optar por definir o bloco como bloco congestionado ou não.
O código padrão poderia definir automaticamente o flag de bloco congestionado
como ``1'' se o mempool estiver acima de algum tamanho e a taxa média desse
conjunto de tamanhos estiver acima de algum valor. Entretanto, um minerador
tem total discricionariedade para alterar as regras de quando algo é
automaticamente definido como bloco congestionado, ou pode optar por definir
permanentemente o flag de congestionamento como permanentemente ligado ou
desligado. Espera-se que a maioria dos mineradores honestos use o
comportamento padrão definido em seus mineradores e não organize um ataque
de 51\%.

Por exemplo, se uma saída de transação pai é gasta por um filho com um valor
de nSequence de 10, deve-se esperar 10 confirmações antes que a transação
se torne válida. Entretanto, se o flag de timestop foi definido, a contagem
de confirmações para, mesmo com novos blocos. Se 6 confirmações se
passaram (mais 4 são necessárias para que a transação seja válida), e o
bloco timestop foi definido no 7º bloco, esse bloco não conta para o
requisito de nSequence de 10 confirmações; o filho ainda está em 6 blocos
para o valor de confirmação relativa. Funcionalmente, isso será armazenado
como algum tipo de altura de bloco auxiliar de timestop usada apenas para
rastrear o valor do timestop. Quando o bit de timestop estiver definido,
todas as transações que usam um valor de nSequence pararão de contar até que
o bit de timestop seja desativado. Isso fornece tempo e espaço de bloco
suficientes para que transações na altura de bloco auxiliar atual de
timestop entrem na blockchain, o que pode evitar que atacantes sistêmicos
ataquem o sistema com sucesso.

Entretanto, isso requer algum tipo de flag no bloco para designar se ele é
um bloco timestop. Para compatibilidade completa com SPV (Simple Payment
Verification; clientes leves), é desejável que isso esteja dentro do
cabeçalho de bloco de 80 bytes, em vez de no coinbase. Existem dois lugares
que podem ser bons para colocar esse flag no cabeçalho do bloco: no horário
do bloco e na versão do bloco. O horário do bloco pode não ser seguro porque
os últimos bits são usados como fonte de entropia por alguns mineradores
ASIC, então um bit pode precisar ser consumido para flags de timestop. Outra
opção seria codificar a ativação de timestop como uma regra rígida de
consenso (por exemplo, via tamanho do bloco), no entanto isso pode tornar
as coisas menos flexíveis. Ao definir padrões sensatos para as regras de
timestop, essas regras podem ser alteradas sem soft-forks de consenso.

Se a versão do bloco for usada como um flag, a informação contextual deve
corresponder ao Chain ID usado em algumas moedas com mineração combinada
(merge-mined).

\subsubsection{Transações de Compromisso Revogáveis}

Ao combinar a atribuição de culpa com a transação revogável, pode-se
determinar quando uma parte não está cumprindo os termos do contrato e
impor penalidades sem confiar na contraparte.

%Diagrama - Transação de Financiamento e todos os gastos do Compromisso 1
\begin{figure}[H]
	\makebox[\linewidth]{
		\includegraphics[width=1\linewidth]{figures/funding-full.pdf}
	}
	\caption{A Transação de Financiamento F, designada em verde, é
		transmitida na blockchain após todas as outras transações
		serem assinadas. Todas as transações que apenas Alice pode
		transmitir estão em roxo. Todas as transações que apenas Bob
		pode transmitir estão em azul. Apenas a Transação de
		Financiamento é transmitida na blockchain neste momento.
	}
\end{figure}

A intenção de criar uma nova Transação de Compromisso é invalidar todas as
Transações de Compromisso antigas ao atualizar o novo saldo com uma nova
Transação de Compromisso. A invalidação de transações antigas pode ocorrer
fazendo de uma saída um Contrato de Maturidade de Sequência Revogável
(RSMC). Para invalidar uma transação, uma transação substitutiva será
assinada e trocada por ambas as partes, que entrega todos os fundos à
contraparte caso uma transação mais antiga seja incorretamente transmitida.
A transmissão incorreta é identificada criando duas Transações de
Compromisso diferentes com as mesmas saídas de saldo final, no entanto, o
pagamento a si mesmo é onerado por um RSMC.

Na prática, existem duas Transações de Compromisso a partir de uma única
saída 2-de-2 da Transação de Financiamento. Dessas duas Transações de
Compromisso, apenas uma pode entrar na blockchain. Cada parte do canal tem
uma versão desse contrato. Assim, se este é o primeiro par de Transações de
Compromisso, a Transação de Compromisso de Alice é definida como C1a, e a
de Bob como C1b. Ao transmitir uma Transação de Compromisso, está-se
solicitando o fechamento e o encerramento do canal. As duas primeiras
saídas da Transação de Compromisso incluem uma Transação de Entrega
(pagamento) do saldo presente não alocado às contrapartes do canal. Se
Alice transmite C1a, uma das saídas pode ser gasta por D1a, que envia
fundos para Bob. Para Bob, C1b pode ser gasta por D1b, que envia fundos a
Alice. A Transação de Entrega (D1a/D1b) é imediatamente resgatável e não é
onerada de forma alguma caso a Transação de Compromisso seja transmitida.

Para a Transação de Compromisso de cada parte, eles atestam que estão
transmitindo a Transação de Compromisso mais recente que possuem. Como
estão atestando que esse é o saldo atual, o saldo pago à contraparte é
assumido como verdadeiro, já que não há benefício direto em pagar fundos à
contraparte como penalidade.

O saldo pago à pessoa que transmitiu a Transação de Compromisso, porém, é
não verificado. Os participantes na blockchain não têm como saber se a
Transação de Compromisso é a mais recente ou não. Se eles não transmitirem
sua versão mais recente, serão penalizados ao terem todos os fundos do
canal entregues à contraparte. Como seus próprios fundos estão onerados em
seu próprio RSMC, eles só poderão reivindicar seus fundos após um número
estipulado de confirmações depois que a Transação de Compromisso for
incluída em um bloco (em nosso exemplo, 1000 confirmações). Se eles
transmitirem sua Transação de Compromisso mais recente, não deve haver
transação de revogação substituindo a transação revogável, então eles
poderão receber seus fundos após algum período definido de tempo (1000
confirmações).

Sabendo quem transmitiu a Transação de Compromisso e onerando seus próprios
pagamentos para que fiquem bloqueados por um período de tempo predefinido,
ambas as partes poderão revogar a Transação de Compromisso no futuro.

\subsubsection{Resgatando Fundos do Canal: Contrapartes Cooperativas}

Qualquer parte pode resgatar os fundos do canal. Entretanto, a parte que
transmite a Transação de Compromisso deve esperar pelo número predefinido de
confirmações descrito no RSMC. A contraparte que não transmitiu a Transação
de Compromisso pode resgatar os fundos imediatamente.

Por exemplo, se a Transação de Financiamento é comprometida com 1 BTC
(metade para cada contraparte) e Bob transmite a Transação de Compromisso
mais recente, C1b, ele deve esperar 1000 confirmações para receber seus
0,5 BTC, enquanto Alice pode gastar 0,5 BTC. Para Alice, essa transação
está totalmente encerrada se Alice concordar que Bob transmitiu a Transação
de Compromisso correta (C1b).

\begin{figure}[H]
	\makebox[\linewidth]{
		\includegraphics[width=1\linewidth]{figures/funding-full-bob-spend.pdf}
	}
	\caption{Quando Bob transmite C1b, Alice pode resgatar sua parte
		imediatamente. Bob deve esperar 1000 confirmações. Quando o
		bloco é imediatamente transmitido, ele está neste estado.
		Transações em verde são transações comprometidas na
		blockchain.
	}
\end{figure}

Após a Transação de Compromisso estar na blockchain por 1000 blocos, Bob
pode então transmitir a transação de Entrega Revogável. Ele deve esperar
1000 blocos para provar que não revogou essa Transação de Compromisso
(C1b). Após 1000 blocos, a transação de Entrega Revogável poderá ser
incluída em um bloco. Se uma parte tentar incluir a transação de Entrega
Revogável em um bloco antes de 1000 confirmações, a transação será
inválida até que se passem 1000 confirmações (quando se tornará válida se
a saída ainda não tiver sido resgatada).

\begin{figure}[H]
	\makebox[\linewidth]{
		\includegraphics[width=1\linewidth]{figures/funding-full-bob-redeem.pdf}
	}
	\caption{Alice concorda que Bob transmitiu a Transação de Compromisso
		correta e 1000 confirmações se passaram. Bob então pode
		transmitir a transação de Entrega Revogável (RD1b) na
		blockchain.
	}
\end{figure}

Depois que Bob transmite a transação de Entrega Revogável, o canal é
totalmente fechado tanto para Alice quanto para Bob; todos receberam os
fundos que ambos concordam ser o saldo atual que cada um possui no canal.

Se em vez disso fosse Alice quem transmitiu a Transação de Compromisso
(C1a), seria ela quem deveria esperar 1000 confirmações em vez de Bob.

\subsubsection{Criando uma nova Transação de Compromisso e Revogando
Compromissos Anteriores}

Embora cada parte possa fechar a Transação de Compromisso mais recente a
qualquer momento, elas também podem optar por criar uma nova Transação de
Compromisso e invalidar a antiga.

Suponha que Alice e Bob queiram agora atualizar seus saldos atuais de 0,5
BTC cada para 0,6 BTC para Bob e 0,4 BTC para Alice. Quando ambos
concordam em fazê-lo, geram um novo par de Transações de Compromisso.

\begin{figure}[H]
	\makebox[\linewidth]{
		\includegraphics[width=1\linewidth]{figures/newcommit-simple.pdf}
	}
	\caption{Quatro transações possíveis podem existir, um par com os
		compromissos antigos e outro par com os novos compromissos.
		Cada parte dentro do canal só pode transmitir metade dos
		compromissos totais (dois para cada). Não há imposição
		explícita impedindo a transmissão de um Compromisso
		específico além dos gastos de penalidade, pois todos são
		gastos não transmitidos válidos. O Compromisso Revogável ainda
		existe com o par C1a/C1b, mas não é exibido por brevidade.
	}
\end{figure}

Quando um novo par de Transações de Compromisso (C2a/C2b) é acordado,
ambas as partes assinam e trocam assinaturas para a nova Transação de
Compromisso e, em seguida, invalidam a antiga. Essa invalidação ocorre
fazendo com que ambas as partes assinem uma Transação de Remediação de
Quebra (Breach Remedy Transaction, BR1), que substitui a Transação de
Entrega Revogável (RD1). Cada parte entrega à outra uma revogação
meio-assinada (BR1) a partir de sua própria Entrega Revogável (RD1), que é
um gasto a partir da Transação de Compromisso. A Transação de Remediação
de Quebra enviará todas as moedas à contraparte no saldo atual do canal.
Por exemplo, se Alice e Bob geram um novo par de Transações de Compromisso
(C2a/C2b) e invalidam compromissos anteriores (C1a/C1b), e posteriormente
Bob transmite incorretamente C1b na blockchain, Alice pode tomar todo o
dinheiro de Bob no canal. Alice pode fazer isso porque Bob provou a Alice
via penalidade que ele nunca transmitirá C1b, pois no momento em que ele
transmitir C1b, Alice é capaz de tomar todo o dinheiro de Bob no canal. Na
prática, ao construir uma transação de Remediação de Quebra para a
contraparte, atestou-se que não se transmitirá nenhum compromisso anterior.
A contraparte pode aceitar isso, porque receberá todo o dinheiro do canal
quando esse acordo for violado.

\begin{figure}[H]
	\makebox[\linewidth]{
		\includegraphics[width=1.3\linewidth]{figures/newcommit-br.pdf}
	}
	\caption{
		Quando C2a e C2b existem, ambas as partes trocam Transações
		de Remediação de Quebra. Ambas as partes agora têm incentivo
		econômico explícito para evitar a transmissão de Transações
		de Compromisso antigas (C1a/C1b). Se qualquer parte desejar
		fechar o canal, usará apenas C2a (Alice) ou C2b (Bob). Se
		Alice transmitir C1a, todo o dinheiro dela vai para Bob. Se
		Bob transmitir C1b, todo o seu dinheiro vai para Alice. Veja
		a figura anterior para as saídas de C2a/C2b.
	}
\end{figure}

Devido a esse fato, é provável que se apague todas as Transações de
Compromisso anteriores quando uma Transação de Remediação de Quebra tiver
sido entregue à contraparte. Se alguém transmitir uma Transação de
Compromisso incorreta (depreciada e invalidada), todo o dinheiro irá para
a contraparte. Por exemplo, se Bob transmitir C1b, contanto que Alice
monitore a blockchain dentro do número predefinido de blocos (neste caso,
1000 blocos), Alice poderá tomar todo o dinheiro deste canal transmitindo
RD1b. Mesmo que o saldo atual do estado do Compromisso (C2a/C2b) seja
0,4 BTC para Alice e 0,6 BTC para Bob, como Bob violou os termos do
contrato, todo o dinheiro vai para Alice como penalidade. Funcionalmente, a
Transação Revogável atua como prova para a blockchain de que Bob violou os
termos do canal, e isso é programaticamente julgado pela blockchain.

\begin{figure}[H]
	\makebox[\linewidth]{
		\includegraphics[width=1.3\linewidth]{figures/newcommit-penalty.pdf}
	}
	\caption{
		Transações em verde estão comprometidas na blockchain. Bob
		transmite incorretamente C1b (apenas Bob pode transmitir
		C1b/C2b). Como ambos concordaram que o estado atual é o par
		de Compromissos C2a/C2b, e atestaram, a cada parte, que
		compromissos antigos foram invalidados por meio de Transações
		de Remediação de Quebra, Alice pode transmitir BR1b e tomar
		todo o dinheiro do canal, contanto que o faça dentro de 1000
		blocos após C1b ser transmitido.
	}
\end{figure}

Entretanto, se Alice não transmitir BR1b dentro de 1000 blocos, Bob pode
conseguir roubar algum dinheiro, já que sua Transação de Entrega Revogável
(RD1b) se torna válida após 1000 blocos. Quando uma Transação de Compromisso
incorreta é transmitida, apenas a Transação de Remediação de Quebra pode
ser transmitida por 1000 blocos (ou qualquer número de confirmações com
que ambas as partes concordem). Após 1000 confirmações de bloco, tanto a
Remediação de Quebra (BR1b) quanto as Transações de Entrega Revogável (RD1b)
podem ser transmitidas a qualquer momento. As transações de Remediação de
Quebra só têm exclusividade dentro desse período predefinido, e qualquer
tempo após isso é funcionalmente uma expiração do prazo prescricional
\textemdash de acordo com o consenso da blockchain do Bitcoin, o tempo
para disputa acabou.

Por essa razão, deve-se monitorar periodicamente a blockchain para verificar
se a contraparte transmitiu uma Transação de Compromisso invalidada, ou
delegar isso a um terceiro. Um terceiro pode ser delegado dando a ele
apenas a transação de Remediação de Quebra. Esse terceiro pode ser
incentivado a monitorar a blockchain transmitindo tal transação em caso de
malícia da contraparte, dando-lhe alguma taxa na saída. Como o terceiro só
pode agir quando a contraparte está agindo maliciosamente, esse terceiro não
tem poder para forçar o fechamento do canal.

\subsubsection{Processo para Criação de Transações de Compromisso Revogáveis}

Para criar Transações de Compromisso revogáveis, é necessária uma
construção adequada do canal desde o início, assinando apenas transações que
possam ser transmitidas a qualquer momento no futuro, ao mesmo tempo em que
se garante que não se perderá devido a contrapartes não cooperativas ou
maliciosas. Isso requer determinar qual chave pública usar para novos
compromissos, já que o uso de SIGHASH\_NOINPUT exige usar chaves únicas para
cada saída RSMC (e HTLC) de Transação de Compromisso. Usamos $P$ para
designar chaves públicas e $K$ para designar a chave privada correspondente
usada para assinar.

Ao gerar a primeira Transação de Compromisso, Alice e Bob concordam em criar
uma saída multisig a partir de uma Transação de Financiamento com uma única
saída $multisig(P_{AliceF}, P_{BobF})$, financiada com 0,5 BTC de Alice e
Bob, totalizando 1 BTC. Essa saída é uma transação Pay to Script
Hash\cite{p2sh}, que exige que Alice e Bob concordem para gastar a partir da
Transação de Financiamento. Eles ainda não tornam a Transação de
Financiamento (F) gastável. Além disso, $P_{AliceF}$ e $P_{BobF}$ são usados
apenas para a Transação de Financiamento; não são usados para mais nada.

Como a transação de Entrega é apenas uma saída P2PKH (endereços bitcoin que
começam com 1) ou uma transação P2SH (comumente reconhecida como endereços
que começam com 3) que as contrapartes designam antecipadamente, isso pode
ser gerado como uma saída de $P_{AliceD}$ e $P_{BobD}$. Para simplificar,
esses endereços de saída permanecerão os mesmos ao longo do canal, já que os
fundos são totalmente controlados pelo destinatário designado depois que a
Transação de Compromisso entra na blockchain. Se desejado, mas não
necessário, ambas as partes podem atualizar e alterar $P_{AliceD}$ e
$P_{BobD}$ para futuras Transações de Compromisso.

Ambas as partes trocam as chaves públicas que pretendem usar para o RSMC
(e o HTLC descrito em seções futuras) para a Transação de Compromisso. Cada
conjunto de Transações de Compromisso usa suas próprias chaves públicas e
nunca são reutilizadas. Ambas as partes podem já conhecer todas as chaves
públicas futuras usando uma construção de Carteira HD BIP 0032\cite{bip32}
ao trocar Chaves Públicas Mestre durante a construção do canal. Se
desejarem gerar um novo par de Transações de Compromisso C2a/C2b, usam
multisig($P_{AliceRSMC2}$, $P_{BobRSMC2}$) para a saída do RSMC.

Após ambas as partes conhecerem os valores de saída das Transações de
Compromisso, ambas criam o par de Transações de Compromisso, por exemplo
C2a/C2b, mas não trocam assinaturas para as Transações de Compromisso.
Ambas assinam a transação de Entrega Revogável (RD2a/RD2b) e trocam as
assinaturas. Bob assina RD1a e a entrega a Alice (usando $K_{BobRSMC2}$),
enquanto Alice assina RD1b e a entrega a Bob (usando $K_{AliceRSMC2}$).

Quando ambas as partes têm a transação de Entrega Revogável, elas trocam
assinaturas para as Transações de Compromisso. Bob assina C1a usando
$K_{BobF}$ e a entrega a Alice, e Alice assina C1b usando $K_{AliceF}$ e a
entrega a Bob.

Neste ponto, tanto a Transação de Compromisso anterior quanto a nova podem
ser transmitidas; tanto C1a/C1b quanto C2a/C2b são válidas. (Note que
Compromissos mais antigos que o Compromisso anterior são invalidados via
penalidades.) Para invalidar C1a e C1b, ambas as partes trocam assinaturas
de Transação de Remediação de Quebra (BR1a/BR1b) para o compromisso
anterior C1a/C1b. Alice envia BR1a a Bob usando $K_{AliceRSMC1}$, e Bob
envia BR1b a Alice usando $K_{BobRSMC1}$. Quando ambas as assinaturas de
Remediação de Quebra forem trocadas, o estado do canal está agora no
Compromisso atual C2a/C2b e os saldos estão agora comprometidos.

Entretanto, em vez de divulgar as assinaturas de BR1a/BR1b, também é
possível apenas divulgar as chaves privadas à contraparte. Isso é mais
eficiente, conforme descrito mais adiante na seção de armazenamento de
chaves. Pode-se divulgar as chaves privadas usadas na própria Transação de
Compromisso. Por exemplo, se Bob deseja invalidar C1b, ele envia suas
chaves privadas usadas em C1b a Alice (ele \textit{NÃO} divulga suas chaves
usadas em C1a, pois isso permitiria roubo de moedas). De forma similar,
Alice divulga a Bob todas as suas chaves privadas das saídas em C1a para
invalidar C1a.

Se Bob transmitir incorretamente C1b, então, como Alice tem todas as chaves
privadas usadas nas saídas de C1b, ela pode tomar o dinheiro. Entretanto,
apenas Bob é capaz de transmitir C1b. Para prevenir esse risco de roubo de
moedas, Bob deve destruir todas as Transações de Compromisso antigas.

\subsection{Fechando um Canal de Forma Cooperativa}

Ambas as partes podem enviar quantos pagamentos desejarem à sua contraparte,
desde que tenham fundos disponíveis no canal, sabendo que, em caso de
divergências, podem transmitir à blockchain o estado atual a qualquer
momento.

Na grande maioria dos casos, todas as saídas da Transação de Financiamento
nunca serão transmitidas na blockchain. Elas estão lá apenas no caso de a
outra parte não ser cooperativa, assim como um contrato raramente é
executado nos tribunais. Uma capacidade comprovada de o contrato ser
imposto de forma determinística é incentivo suficiente para que ambas as
partes ajam honestamente.

Quando qualquer parte deseja fechar um canal de forma cooperativa, ela
poderá fazê-lo contatando a outra parte e gastando a partir da Transação
de Financiamento com uma saída da Transação de Compromisso mais atual,
diretamente, sem qualquer condição onerando o script. Nenhum pagamento
adicional pode ocorrer no canal.

\begin{figure}[H]
	\makebox[\linewidth]{
		\includegraphics[width=1\linewidth]{figures/cooperative-close.pdf}
	}
	\caption{Se ambas as contrapartes são cooperativas, elas tomam os
		saldos da Transação de Compromisso atual e gastam a partir da
		Transação de Financiamento com uma Transação de Liquidação por
		Exercício (Exercise Settlement, ES). Se em vez disso a
		Transação de Compromisso mais recente for transmitida, o
		pagamento (menos taxas) será o mesmo.
	}
\end{figure}

O propósito do fechamento cooperativo é reduzir o número de transações que
ocorrem na blockchain, e ambas as partes poderão receber seus fundos
imediatamente (em vez de uma parte esperar que a transação de Entrega
Revogável se torne válida).

Os canais podem permanecer em perpetuidade até que decidam encerrar a
transação de forma cooperativa, ou quando uma parte não coopera com a outra
e o canal é fechado e imposto na blockchain.

\subsection{Implicações e Resumo dos Canais Bidirecionais}

Ao garantir que os canais só possam ser atualizados com o consentimento de
ambas as partes, é possível construir canais que existem perpetuamente na
blockchain. Ambas as partes podem atualizar o saldo dentro do canal com
quaisquer saldos de saída desejados, desde que sejam iguais ou inferiores
ao total de fundos comprometidos na Transação de Financiamento; os saldos
podem se mover em ambas as direções. Se uma parte se tornar maliciosa,
qualquer parte pode fechar imediatamente o canal e transmitir o estado mais
atual à blockchain. Usando uma construção de caução de fidelidade
(Transações de Entrega Revogáveis), se uma parte violar os termos do
canal, os fundos serão enviados à contraparte, contanto que a prova da
violação (Transação de Remediação de Quebra) seja inserida na blockchain de
forma tempestiva. Se ambas as partes forem cooperativas, o canal pode
permanecer aberto indefinidamente, possivelmente por muitos anos.

Esse tipo de construção só é possível porque a adjudicação ocorre
programaticamente na blockchain como parte do consenso do Bitcoin, de modo
que não é necessário confiar na outra parte. Como resultado, a contraparte
do canal não detém a custódia ou o controle total dos fundos.

\section{Contrato Timelock com Hash (HTLC)}

Um canal de pagamento bidirecional só permite a transferência segura de
fundos dentro de um canal. Para poder construir transferências seguras
usando uma rede de canais ao longo de múltiplos saltos até o destino final,
é necessária uma construção adicional: o Contrato Timelock com Hash
(Hashed Timelock Contract, HTLC).

O propósito de um HTLC é permitir um estado global entre múltiplos nós via
hashes. Esse estado global é garantido por compromissos de tempo e pelo
desbloqueio temporizado de recursos via divulgação de pré-imagens. O
``travamento'' transacional ocorre globalmente via compromissos; em qualquer
momento no tempo, um único participante é responsável por divulgar ao
próximo participante se tem conhecimento da pré-imagem $R$. Essa construção
não requer confiança custodial na contraparte do canal, nem em qualquer
outro participante da rede.

Para alcançar isso, um HTLC deve ser capaz de criar certas transações que
só são válidas após uma certa data, usando nLockTime, bem como a divulgação
de informações à contraparte do canal. Adicionalmente, esses dados devem ser
revogáveis, pois é necessário poder desfazer um HTLC.

Um HTLC também é um contrato de canal com a contraparte que é exigível via
blockchain. As contrapartes em um canal concordam com os seguintes termos
para um Contrato Timelock com Hash:

\begin{enumerate}
	\item Se Bob conseguir produzir a Alice um dado aleatório de entrada de
		20 bytes desconhecido $R$ a partir de um hash conhecido $H$,
		dentro de três dias, então Alice liquidará o contrato pagando
		a Bob 0,1 BTC.

	\item Se três dias se passarem, a cláusula acima é nula e sem efeito, e
		o processo de compensação é invalidado; ambas as partes não
		devem tentar liquidar e reivindicar pagamento após três dias.

	\item Qualquer parte pode (e deve) pagar conforme os termos deste
		contrato em qualquer método escolhido pelos participantes e
		encerrar este contrato antecipadamente, contanto que ambos os
		participantes concordem.

	\item A violação dos termos acima incorrerá em uma penalidade máxima dos
		fundos bloqueados neste contrato, a serem pagos à contraparte
		não violadora como uma caução de fidelidade.

\end{enumerate}

Para clareza nos exemplos, usamos dias para HTLCs e altura de bloco para
RSMCs. Na realidade, o HTLC também deveria ser definido como altura de bloco
(por exemplo, 3 dias equivalem a 432 blocos).

Na prática, deseja-se construir um pagamento contingente ao conhecimento de
$R$ pelo destinatário dentro de um determinado prazo. Após esse prazo, os
fundos são reembolsados ao remetente.

De forma similar aos RSMCs, esses termos contratuais são impostos
programaticamente na blockchain do Bitcoin e não requerem confiança na
contraparte para aderir aos termos do contrato, já que todas as violações
são penalizadas por meio de cauções de fidelidade unilateralmente
exigíveis, que são construídas usando transações de penalidade que gastam a
partir dos estados de compromisso. Se Bob conhecer $R$ dentro de três dias,
ele pode resgatar os fundos transmitindo uma transação; Alice não tem como
reter os fundos de nenhuma forma, porque o script retorna como válido quando
a transação é gasta na blockchain do Bitcoin.

Um HTLC é uma saída adicional em uma Transação de Compromisso com um script
de saída único:

\begin{lstlisting}
OP_IF
	OP_HASH160 <Hash160(R)> OP_EQUALVERIFY
	2 <Alice2> <Bob2> OP_CHECKMULTISIG
OP_ELSE
	2 <Alice1> <Bob1> OP_CHECKMULTISIG
OP_ENDIF
\end{lstlisting}

Conceitualmente, esse script tem dois caminhos possíveis de gasto a partir
de uma única saída HTLC. O primeiro caminho (definido no OP\_IF) envia
fundos a Bob se Bob produzir $R$. O segundo caminho é resgatado usando um
reembolso a Alice com timelock de 3 dias. O timelock de 3 dias é imposto
usando nLockTime na transação que gasta.

\subsection{Construção de HTLC Não Revogável}

\begin{figure}[H]
	\makebox[\linewidth]{
		\includegraphics[width=1\linewidth]{figures/htlc-concept.pdf}
	}
	\caption{Esta é uma implementação ingênua e não funcional de um HTLC.
		Apenas o caminho do HTLC a partir da Transação de Compromisso
		é exibido. Note que há dois gastos possíveis a partir de uma
		saída HTLC. Se Bob conseguir produzir a pré-imagem $R$ dentro
		de 3 dias, ele pode resgatar o caminho 1. Após três dias,
		Alice pode transmitir o caminho 2. Quando 3 dias se passam,
		qualquer um é válido. Esse modelo, contudo, não funciona com
		múltiplas Transações de Compromisso.
	}
\end{figure}

Se $R$ for produzido dentro de 3 dias, Bob pode resgatar os fundos
transmitindo a transação ``Entrega''. Um requisito para que a transação
``Entrega'' seja válida exige que $R$ seja incluído com a transação. Se $R$
não for incluído, a transação ``Entrega'' é inválida. No entanto, se 3 dias
se passaram, os fundos podem ser devolvidos a Alice transmitindo a transação
``Timeout''. Quando 3 dias se passaram e $R$ foi divulgado, qualquer das
transações pode ser válida.

É responsabilidade individual de ambas as partes garantir que possam fazer
sua transação entrar na blockchain para garantir que os saldos estejam
corretos. Para Bob, a fim de receber os fundos, ele deve transmitir a
transação ``Entrega'' na blockchain do Bitcoin, ou então liquidar com Alice
(cancelando o HTLC). Para Alice, ela deve transmitir o ``Timeout'' daqui a
3 dias para receber o reembolso, ou cancelar o HTLC inteiramente com Bob.

Ainda assim, esse tipo de construção simplista tem problemas semelhantes
aos de uma construção incorreta de canal de pagamento bidirecional. Quando
uma Transação de Compromisso antiga é transmitida, qualquer parte pode
tentar roubar fundos, já que ambos os caminhos podem ser válidos a
posteriori. Por exemplo, se $R$ for divulgado 1 ano depois e uma Transação
de Compromisso incorreta for transmitida, ambos os caminhos são válidos e
resgatáveis por qualquer das partes; o contrato ainda não é exigível na
blockchain. O encerramento do HTLC é absolutamente necessário, porque para
Alice obter seu reembolso, ela deve encerrar o contrato e recebê-lo. Caso
contrário, quando Bob descobrir $R$ depois que 3 dias se passaram, ele
poderá roubar os fundos que deveriam ir para Alice. Com contrapartes não
cooperativas, não é possível encerrar um HTLC sem transmiti-lo à blockchain
do Bitcoin, já que a parte não cooperativa não estará disposta a criar uma
nova Transação de Compromisso.

\subsection{HTLC Revogável Off-chain}

Para poder encerrar esse contrato off-chain sem uma transmissão à
blockchain do Bitcoin, é necessário embutir RSMCs na saída, que terão uma
construção similar à do canal bidirecional.

\begin{figure}[H]
	\vspace*{-2cm}
	\makebox[\linewidth]{
		\includegraphics[width=1.3\linewidth]{figures/htlc-functional.pdf}
	}
	\caption{Se Alice transmitir C2a, então a metade esquerda será
		executada. Se Bob transmitir C2b, então a metade direita
		será executada. Qualquer parte pode transmitir sua Transação
		de Compromisso a qualquer momento. O Timeout do HTLC só é
		válido após 3 dias. As Execuções do HTLC só podem ser
		transmitidas se a pré-imagem do hash $R$ for conhecida.
		Compromissos anteriores (e suas transações dependentes) não
		são exibidos por brevidade.
	}
\end{figure}

Suponha que Alice e Bob queiram atualizar seu saldo no canal no
Compromisso 1, com um saldo de 0,5 para Alice e 0,5 para Bob.

Alice deseja enviar 0,1 a Bob contingente ao conhecimento de $R$ dentro de
3 dias; após 3 dias, ela quer seu dinheiro de volta caso Bob não produza
$R$.

A nova Transação de Compromisso terá um reembolso total do saldo atual a
Alice e Bob (Saídas 0 e 1), com a saída 2 sendo o HTLC, que descreve os
fundos em trânsito. Como 0,1 estará onerado em um HTLC, o saldo de Alice é
reduzido para 0,4 e o de Bob permanece o mesmo em 0,5.

Essa nova Transação de Compromisso (C2a/C2b) terá uma saída HTLC com dois
gastos possíveis. Cada gasto é diferente, dependendo da versão da Transação
de Compromisso de cada contraparte. De forma similar ao canal de
pagamento bidirecional, quando uma parte transmite seu Compromisso, os
pagamentos à contraparte serão presumidos válidos e não serão invalidados.
Isso pode ocorrer porque, ao transmitir uma Transação de Compromisso,
está-se atestando que essa é a Transação de Compromisso mais recente. Se
for a mais recente, está-se também atestando que o HTLC existe e não foi
invalidado antes; portanto, pagamentos potenciais à contraparte devem ser
válidos.

Note que os nomes de transações HTLC (começando com a letra H) começarão
com o número 1, cujos valores não se correlacionam com as Transações de
Compromisso. Esta é simplesmente a primeira transação HTLC. As transações
HTLC podem persistir entre Transações de Compromisso. Cada HTLC tem 4
chaves por lado da transação (C2a e C2b), totalizando 8 chaves por
contraparte.

A saída HTLC na Transação de Compromisso tem dois conjuntos de chaves por
contraparte na saída.

Para a Transação de Compromisso de Alice (C2a), o script de saída HTLC
requer $multisig(P_{Alice2}, P_{Bob2})$ onerado pela divulgação de $R$,
bem como $multisig(P_{Alice1}, P_{Bob1})$ sem oneração.

Para a Transação de Compromisso de Bob (C2b), o script de saída HTLC requer
$multisig(P_{Alice6}, P_{Bob6})$ onerado pela divulgação de $R$, bem como
$multisig(P_{Alice5}, P_{Bob5})$ sem oneração.

Os estados de saída do HTLC são diferentes dependendo de qual Transação de
Compromisso é transmitida.

\subsubsection{HTLC quando o Remetente Transmite a Transação de Compromisso}

Para o remetente (Alice), a transação ``Entrega'' é enviada como uma
transação de Entrega de Execução do HTLC (HTLC Execution Delivery, HED1a),
que não é onerada em um RSMC. Assume-se que esse HTLC nunca foi encerrado
off-chain, já que Alice atesta que a Transação de Compromisso transmitida é
a mais recente. Se Bob conseguir produzir a pré-imagem $R$, ele poderá
resgatar fundos do HTLC depois que a Transação de Compromisso for
transmitida na blockchain. Esta transação consome $multisig(P_{Alice2},
P_{Bob2})$ se Alice transmitir seu Compromisso C2a. Apenas Bob pode
transmitir HED1a, já que apenas Alice deu sua assinatura de HED1a a Bob.

Entretanto, se 3 dias se passaram desde a formação do HTLC, Alice poderá
transmitir uma transação ``Timeout'', a transação de Timeout do HTLC
(HT1a). Essa transação é um RSMC. Consome a saída $multisig(P_{Alice1},
P_{Bob1})$ sem exigir divulgação de $R$ se Alice transmitir C2a. Essa
transação não pode entrar na blockchain até que 3 dias se passem. A saída
para essa transação é um RSMC com $multisig(P_{Alice3}, P_{Bob3})$ com
maturidade relativa de 1000 blocos, e $multisig(P_{Alice4}, P_{Bob4})$ sem
requisito de maturidade de confirmação. Apenas Alice pode transmitir HT1a,
já que apenas Bob deu sua assinatura de HT1a a Alice.

Após HT1a entrar na blockchain e 1000 confirmações de bloco ocorrerem, uma
transação de Entrega Revogável do Timeout do HTLC (HTLC Timeout Revocable
Delivery, HTRD1a) pode ser transmitida por Alice, consumindo
$multisig(P_{Alice3}, P_{Bob3})$. Apenas Alice pode transmitir HTRD1a 1000
blocos após HT1a ser transmitida, já que apenas Bob deu sua assinatura de
HTRD1a a Alice. Essa transação pode ser revogável quando outra transação
substitui HTRD1a usando $multisig(P_{Alice4}, P_{Bob4})$, que não tem
nenhum requisito de maturidade de bloco.

\subsubsection{HTLC quando o Destinatário Transmite a Transação de Compromisso}

Para o destinatário potencial (Bob), o ``Timeout'' de recebimento é
reembolsado como uma transação de Entrega de Timeout do HTLC (HTLC Timeout
Delivery, HTD1b). Essa transação reembolsa diretamente os fundos ao
remetente original (Alice) e não é onerada em um RSMC. Assume-se que esse
HTLC nunca foi encerrado off-chain, já que Bob atesta que a Transação de
Compromisso transmitida (C2b) é a mais recente. Se 3 dias se passaram,
Alice pode transmitir HTD1b e levar o reembolso. Essa transação consome
$multisig(P_{Alice5}, P_{Alice5})$ se Bob transmitir C2b. Apenas Alice pode
transmitir HTD1b, já que Bob deu sua assinatura de HTD1b a Alice.

Entretanto, se HTD1b não for transmitida (3 dias não se passaram) e Bob
conhecer a pré-imagem $R$, então Bob poderá transmitir a transação de
Execução do HTLC (HE1b) se conseguir produzir $R$. Essa transação é um
RSMC. Consome a saída $multisig(P_{Alice6}, P_{Bob6})$ e requer divulgação
de $R$ se Bob transmitir C2b. A saída dessa transação é um RSMC com
$multisig(P_{Alice7}, P_{Bob7})$ com maturidade relativa de 1000 blocos, e
$multisig(P_{Alice8}, P_{Bob8})$ sem requisitos de maturidade de bloco.
Apenas Bob pode transmitir HE1b, já que apenas Alice deu sua assinatura de
HE1b a Bob.

Após HE1b entrar na blockchain e 1000 confirmações de bloco ocorrerem,
uma transação de Entrega Revogável de Execução do HTLC (HTLC Execution
Revocable Delivery, HERD1b) pode ser transmitida por Bob, consumindo
$multisig(P_{Alice7}, P_{Bob7})$. Apenas Bob pode transmitir HERD1b 1000
blocos após HE1b ser transmitida, já que apenas Alice deu sua assinatura de
HERD1b a Bob. Essa transação pode ser revogável quando outra transação
substitui HERD1b usando $multisig(P_{Alice8}, P_{Bob8})$, que não tem
requisitos de maturidade de bloco.

\subsection{Encerramento Off-chain do HTLC}

Depois que um HTLC é construído, encerrar um HTLC off-chain requer que ambas
as partes concordem sobre o estado do canal. Se o destinatário puder provar
conhecimento de $R$ à contraparte, o destinatário está provando que é capaz
de fechar imediatamente o canal na blockchain do Bitcoin e receber os
fundos. Nesse ponto, se ambas as partes desejarem manter o canal aberto,
elas devem encerrar o HTLC off-chain e criar uma nova Transação de
Compromisso refletindo o novo saldo.

\begin{figure}[H]
	%\vspace*{-2cm}
	\makebox[\linewidth]{
		\includegraphics[width=1.2\linewidth]{figures/htlc-newcommit.pdf}
	}
	\caption{Como Bob provou a Alice que conhece $R$ ao dizer-lhe $R$,
		Alice está disposta a atualizar o saldo com uma nova
		Transação de Compromisso. O pagamento será o mesmo, seja C2
		ou C3 transmitida neste momento.
	}
\end{figure}

De forma similar, se o destinatário não puder provar conhecimento de $R$
divulgando $R$, ambas as partes devem concordar em encerrar o HTLC e criar
uma nova Transação de Compromisso com o saldo do HTLC reembolsado ao
remetente.

Se as contrapartes não conseguirem chegar a um acordo ou se tornarem de
outra forma não responsivas, elas devem fechar o canal transmitindo as
transações necessárias do canal na blockchain do Bitcoin.

Entretanto, se forem cooperativas, podem fazer isso primeiro gerando uma
nova Transação de Compromisso com os novos saldos e, em seguida,
invalidando o Compromisso anterior trocando transações de Remediação de
Quebra (BR2a/BR2b). Adicionalmente, se estiverem encerrando um HTLC
específico, também devem trocar algumas de suas próprias chaves privadas
usadas nas transações do HTLC.

Por exemplo, se Alice deseja encerrar o HTLC, ela divulgará $K_{Alice1}$ e
$K_{Alice4}$ a Bob. Correspondentemente, se Bob deseja encerrar o HTLC,
ele divulgará $K_{Bob6}$ e $K_{Bob8}$ a Alice. Após as chaves privadas
serem divulgadas à contraparte, se Alice transmitir C2a, Bob poderá tomar
todos os fundos do HTLC imediatamente. Se Bob transmitir C2b, Alice poderá
tomar todos os fundos do HTLC imediatamente. Note que, quando um HTLC é
encerrado, a Transação de Compromisso mais antiga também deve ser
revogada.

\begin{figure}[H]
	\vspace*{-3cm}
	\makebox[\linewidth]{
		\includegraphics[width=1.3\linewidth]{figures/Diagram2.pdf}
	}
	\caption{Uma Transação de Compromisso totalmente revogada e um HTLC
		encerrado. Se qualquer parte transmitir o Compromisso 2,
		perderá todo seu dinheiro para a contraparte. Outros
		compromissos (por exemplo, se o Compromisso 3 for o atual)
		não são exibidos por brevidade.
	}
\end{figure}

Como ambas as partes podem provar o estado atual uma à outra, elas podem
chegar a um acordo sobre o saldo atual dentro do canal. Como podem
transmitir o estado atual na blockchain, são capazes de chegar a um acordo
sobre a liquidação líquida e o encerramento do HTLC com uma nova Transação
de Compromisso.

\subsection{Formação e Ordem de Fechamento do HTLC}

Para criar um novo HTLC, o processo é o mesmo de criar uma nova Transação
de Compromisso, exceto que as assinaturas do HTLC são trocadas antes das
assinaturas da nova Transação de Compromisso.

Para fechar um HTLC, o processo é o seguinte (de C2 para C3):

\begin{enumerate}
	\item Alice assina e envia sua assinatura para RD3b e C3b. Nesse
		ponto Bob pode optar por transmitir C3b ou C2b (com o HTLC)
		com o mesmo pagamento. Bob está disposto, após receber C3b,
		a encerrar C2b.
	\item Bob assina e envia sua assinatura para RD3a e C3a, bem como suas
		chaves privadas usadas para o Compromisso 2 e o HTLC sendo
		encerrado; ele envia a Alice $K_{BobRSMC2}$, $K_{Bob5}$ e
		$K_{Bob8}$. Nesse ponto, Bob só deve transmitir C3b e não
		deve transmitir C2b, pois perderá todo seu dinheiro se o
		fizer. Bob revogou totalmente C2b e o HTLC. Alice está
		disposta, após receber C3a, a encerrar C2b.
	\item Alice assina e envia sua assinatura para RD3b e C3b, bem como
		suas chaves privadas usadas para o Compromisso 2 e o HTLC
		sendo encerrado; ela envia a Bob $K_{AliceRSMC2}$, $K_{Bob1}$
		e $K_{Bob4}$. Nesse ponto, nenhuma das partes deve transmitir
		o Compromisso 2; se o fizerem, seus fundos irão para a
		contraparte. O antigo Compromisso e o antigo HTLC estão agora
		revogados e totalmente encerrados. Apenas o novo Compromisso
		3 permanece, e ele não tem um HTLC.
\end{enumerate}

Quando o HTLC tiver sido fechado, os fundos são atualizados de modo que o
saldo atual no canal seja aquele que existiria se o contrato HTLC tivesse
sido concluído e transmitido na blockchain. Em vez disso, ambas as partes
optam por fazer a novação off-chain e atualizar seus pagamentos dentro do
canal.

É absolutamente necessário que ambas as partes concluam a novação off-chain
dentro de sua janela de tempo designada. Para o destinatário (Bob), ele
deve conhecer $R$ e atualizar seu saldo com Alice dentro de 3 dias (ou
qualquer tempo selecionado); caso contrário, Alice poderá resgatá-lo
dentro de 3 dias. Para Alice, muito logo depois que seu timeout se tornar
válido, ela deve novar ou transmitir a transação de Timeout do HTLC. Ela
também deve novar ou transmitir a transação de Entrega Revogável de
Timeout do HTLC assim que se tornar válida. Se a contraparte não estiver
disposta a novar ou estiver protelando, então é necessário transmitir o
estado atual do canal (incluindo as transações HTLC) à blockchain do
Bitcoin.

A flexibilidade de tempo dessas ofertas para novar depende das dependências
contingentes do hashlock $R$. Se for estabelecido um contrato em que o HTLC
deva ser resolvido dentro de 1 dia, então, se a transação expirar, Alice
deve resolvê-la até o dia 4 (3 dias mais 1), caso contrário, Alice corre
risco de perder fundos.

\section{Armazenamento de Chaves}

As chaves são geradas usando Carteiras Hierárquicas Determinísticas BIP
0032\cite{bip32}. As chaves são pré-geradas por ambas as partes. As chaves
são geradas em uma árvore de Merkle e estão muito profundamente na árvore.
Por exemplo, Alice pré-gera um milhão de chaves, cada chave sendo um filho
da chave anterior. Alice aloca quais chaves usar de acordo com algum modo
determinístico. Por exemplo, ela começa com o filho mais profundo da árvore
para gerar muitas subchaves para o dia 1. Essa chave é usada como uma chave
mestre para todas as chaves geradas no dia 1. Ela entrega a Bob o endereço
que deseja usar para a próxima transação, e divulga a chave privada a Bob
quando ela for invalidada. Quando Alice divulga a Bob todas as chaves
privadas derivadas da chave mestre do dia 1 e não deseja mais continuar
usando essa chave mestre, ela pode divulgar a chave mestre do dia 1 a Bob.
Nesse ponto, Bob não precisa armazenar todas as chaves derivadas da chave
mestre do dia 1. Bob faz o mesmo para Alice e lhe entrega sua chave do
dia 1.

Quando todas as chaves privadas do Dia 2 tiverem sido trocadas, por
exemplo, no dia 5, Alice divulga sua chave do Dia 2. Bob é capaz de gerar
a chave do Dia 1 a partir da chave do Dia 2, já que a chave do Dia 1
também é um filho da chave do Dia 2.

Se uma contraparte transmitir a Transação de Compromisso errada, qual chave
privada usar em uma transação para recuperar fundos pode ser obtida por
força bruta ou, se ambas as partes concordarem, podem usar o número de id
de sequência ao criar a transação para identificar quais conjuntos de
chaves são usados.

Isso permite que os participantes de um canal tenham estados de saída
(transações) anteriores invalidados por ambas as partes sem usar muitos
dados. Ao divulgar chaves privadas pré-organizadas em uma árvore de
Merkle, é possível invalidar milhões de transações antigas com apenas
alguns quilobytes de dados por canal. Canais centrais na Lightning Network
podem conduzir bilhões de transações sem necessidade de custos
significativos de armazenamento.

\section{Taxas de Transação da Blockchain para Canais Bidirecionais}

É possível para cada participante gerar versões diferentes de transações
para atribuir culpa sobre quem transmitiu a transação na blockchain. Tendo
conhecimento de quem transmitiu uma transação e a capacidade de atribuir
culpa, um serviço de terceiros pode ser usado para reter taxas em um
escrow multisig 2-de-3. Se alguém deseja transmitir a cadeia de transações
em vez de concordar em fazer um Fechamento de Financiamento ou
substituição por uma nova Transação de Compromisso, comunica-se com o
terceiro e transmite a cadeia à blockchain. Se a contraparte recusar o
aviso do terceiro para cooperar, a penalidade é recompensada à parte não
cooperativa. Na maioria dos casos, os participantes podem ser indiferentes
às taxas de transação no caso de uma contraparte não cooperativa.

Deve-se escolher contrapartes no canal que sejam cooperativas, mas isso
não é uma necessidade absoluta para o funcionamento do sistema. Note que
isso não exige confiança no restante da rede e só é relevante para as
relativamente pequenas taxas de transação. A parte menos confiável pode
ser justamente a responsável pelas taxas de transação.

As taxas da Lightning Network provavelmente serão significativamente
menores do que as taxas de transação da blockchain. As taxas são
amplamente derivadas do valor temporal do bloqueio de fundos para uma rota
específica, bem como do pagamento pela chance de fechamento do canal na
blockchain. Estas devem ser significativamente menores que as transações
on-chain, já que muitas transações em um canal da Lightning Network podem
ser liquidadas em uma única transação na blockchain. Com uma rede
suficientemente robusta e interconectada, as taxas devem
assintoticamente se aproximar de níveis desprezíveis para muitos tipos de
transação. Com taxas baixas e transações rápidas, será possível construir
micropagamentos escaláveis, mesmo entre sistemas de alta frequência, como
aplicações de Internet das Coisas ou microcobrança por unidade.

\section{Pagar ao Contrato (Pay to Contract)}

É possível construir um contrato de ``Entrega contra Pagamento''
criptograficamente comprovável, ou pay-to-contract\cite{paytocontract},
como prova de pagamento. Essa prova pode ser estabelecida como o
conhecimento da entrada R a partir de hash(R) como pagamento de um certo
valor. Ao embutir uma cláusula no contrato entre comprador e vendedor
declarando que conhecer R é prova de fundos enviados, o destinatário dos
fundos não tem incentivo para divulgar R, a menos que tenha certeza de que
receberá o pagamento. Quando os fundos forem eventualmente puxados do
comprador pela sua contraparte no canal de micropagamento, R é divulgado
como parte dessa retirada de fundos. Pode-se elaborar documentos legais
em papel que especifiquem que o conhecimento ou a divulgação de R implica
o cumprimento do pagamento. O remetente pode então arranjar um contrato
criptograficamente assinado com conhecimento de entradas para hashes
tratado como cumprimento do contrato em papel antes que o pagamento
ocorra.

\section{A Lightning Network do Bitcoin}

Tendo um canal de micropagamento com contratos onerados por hashlocks e
timelocks, é possível compensar transações em uma rede de pagamentos com
múltiplos saltos usando uma série de timelocks decrescentes, sem câmaras
de compensação centrais adicionais.

Tradicionalmente, os mercados financeiros compensam transações
transferindo a obrigação de entrega a um ponto central e liquidam
transferindo a propriedade por meio desse hub central. Sistemas de
transferência bancária e de fundos (como ACH e a rede de cartões Visa) ou
câmaras de compensação de ações (como a DTCC) operam dessa forma.

Como o Bitcoin permite dinheiro programático, é possível criar transações
sem contatar uma câmara de compensação central. As transações podem ser
executadas off-chain, sem terceiros que coletam todos os fundos antes de
desembolsá-los \textemdash apenas transações com contrapartes de canal não
cooperativas tornam-se automaticamente julgadas na blockchain.

A obrigação de entregar fundos a um destinatário final é obtida por meio
de um processo de delegação encadeada. Cada participante ao longo do
caminho assume a obrigação de entregar a um destinatário específico. Cada
participante passa essa obrigação ao próximo participante no caminho. A
obrigação de cada participante subsequente no caminho, definida em seus
respectivos HTLCs, tem tempo menor para conclusão em comparação com o
participante anterior. Dessa forma, cada participante tem certeza de que
poderá reivindicar os fundos quando a obrigação for enviada ao longo do
caminho.

A Linguagem de Scripts de Transação do Bitcoin, uma forma do que alguns
chamam de implementação de ``Contratos Inteligentes''\cite{smartcontracts},
permite sistemas sem câmaras de compensação custodiais confiáveis ou
serviços de escrow.

\subsection{Timelocks Decrescentes}

Suponha que Alice deseje enviar 0,001 BTC a Dave. Ela localiza uma rota por
Bob e Carol. O caminho de transferência seria Alice para Bob para Carol
para Dave.

\begin{figure}[H]
	\makebox[\linewidth]{
		\includegraphics[width=1\linewidth]{figures/network-fig0.pdf}
	}
	\caption{Pagamento pela Lightning Network usando HTLCs.
	}
\end{figure}

Quando Alice envia pagamento a Dave por Bob e Carol, ela solicita a Dave
hash(R) para usar nesse pagamento. Alice então conta a quantidade de
saltos até o destinatário e usa isso como expiração do HTLC. Neste caso,
ela define a expiração do HTLC em 3 dias. Bob então cria um HTLC com Carol
com expiração de 2 dias, e Carol faz o mesmo com Dave com expiração de 1
dia. Dave agora está livre para divulgar R a Carol, e ambas as partes
provavelmente concordarão em liquidação imediata via novação com uma
Transação de Compromisso substitutiva. Isso então ocorre passo a passo de
volta a Alice. Note que isso ocorre off-chain e nada é transmitido à
blockchain quando todas as partes são cooperativas.

\begin{figure}[H]
	\makebox[\linewidth]{
		\includegraphics[width=1\linewidth]{figures/network-fig1.pdf}
	}
	\caption{Liquidação do HTLC, os fundos de Alice são enviados a Dave.
	}
\end{figure}

Timelocks decrescentes são usados para que todas as partes ao longo do
caminho saibam que a divulgação de R permitirá à parte que divulga puxar
os fundos, já que, no pior caso, estarão puxando fundos após a data até a
qual devem receber R. Se Dave não produzir R dentro de 1 dia para Carol,
então Carol poderá fechar o HTLC. Se Dave transmitir R após 1 dia, ele não
poderá puxar fundos de Carol. A responsabilidade de Carol com Bob ocorre
no dia 2, então Carol nunca será responsável pelo pagamento a Dave sem ser
capaz de puxar fundos de Bob, contanto que atualize sua transação com Dave
via transmissão à blockchain ou via novação.

Caso R seja divulgado aos participantes na metade da expiração ao longo do
caminho (por exemplo, dia 2), é possível que algumas partes ao longo do
caminho sejam enriquecidas. O remetente poderá conhecer R; assim, devido
ao Pay to Contract, o pagamento terá sido cumprido mesmo que o
destinatário não tenha recebido os fundos. Portanto, o destinatário nunca
deve divulgar R a menos que tenha recebido um HTLC de sua contraparte de
canal; eles têm garantia de receber pagamento de uma de suas contrapartes
de canal mediante a divulgação da pré-imagem.

Caso uma parte desconecte completamente, a contraparte será responsável
por transmitir o estado atual da Transação de Compromisso do canal à
blockchain. Apenas o estado do canal com falha (não responsivo) é fechado
na blockchain; todos os outros canais devem continuar a atualizar suas
Transações de Compromisso via novação dentro do canal. Portanto, o risco
de contraparte para taxas de transação é exposto apenas às contrapartes
diretas do canal. Se um nó ao longo do caminho decide se tornar não
responsivo, os participantes não diretamente conectados a esse nó sofrem
apenas uma redução do valor temporal de seus fundos por não realizarem a
liquidação antecipada antes do fechamento do HTLC.

\begin{figure}[H]
	\makebox[\linewidth]{
		\includegraphics[width=1\linewidth]{figures/network-fig2.pdf}
	}
	\caption{Apenas os canais não responsivos são transmitidos na
		blockchain; todos os outros são liquidados off-chain via
		novação.
	}
\end{figure}

\subsection{Valor do Pagamento}

É preferível usar um pagamento pequeno por HTLC. Não se deve usar um
pagamento extremamente alto, caso o pagamento não seja totalmente
roteado até seu destino. Se o pagamento não chegar ao destino e um dos
participantes ao longo do caminho não for cooperativo, é possível que o
remetente precise esperar até a expiração antes de receber um reembolso.
A entrega pode ter perdas, similar a pacotes na internet, mas a rede não
pode roubar fundos em trânsito. Como as transações não atingem a
blockchain com contrapartes de canal cooperativas, recomenda-se usar
pagamentos tão pequenos quanto possível. Existe um trade-off entre
bloquear taxas de transação a cada salto e o desejo de usar um valor de
transação tão pequeno quanto possível (este último pode incorrer em taxas
totais maiores). Transferências menores com mais intermediários implicam
em um percentual maior pago como taxas da Lightning Network aos
intermediários.

\subsection{Falha de Compensação e Reroteamento}

Se uma transação falhar em chegar ao seu destino final, o destinatário deve
enviar um pagamento igual ao remetente com o mesmo hash, mas sem divulgar
R. Isso compensará a divulgação do hash para o remetente, mas pode
não fazê-lo para o destinatário. O destinatário, que gerou o hash, deve
descartar R e nunca transmiti-lo. Se um canal ao longo do caminho não puder
ser contatado, então os canais podem optar por esperar até que o caminho
expire, situação em que todos os participantes provavelmente fecharão o
HTLC como não liquidado, sem qualquer pagamento, com uma nova Transação
de Compromisso.

\begin{figure}[H]
	\makebox[\linewidth]{
		\includegraphics[width=1\linewidth]{figures/network-fig3.pdf}
	}
	\caption{Dave cria um caminho de volta a Alice depois que Alice falha
		em enviar fundos a Dave, porque Carol não é cooperativa. A
		entrada R de hash(R) nunca é transmitida por Dave, porque
		Carol não completou suas ações. Se R fosse transmitido,
		Alice ficaria em equilíbrio (break-even). Dave, que controla
		R, nunca deve transmitir R, porque pode não receber fundos
		de Carol; ele deve deixar os contratos expirarem. Alice e
		Bob também têm a opção de compensar e fechar o contrato
		antecipadamente, neste diagrama.
	}
\end{figure}

Se a rota de reembolso for a mesma que a rota de pagamento, e não houver
contratos meio-assinados em que uma parte possa roubar fundos, é possível
cancelar a transação totalmente substituindo-a por uma nova Transação de
Compromisso, começando pelo nó mais recente que participou do HTLC.

Também é possível liberar um canal criando uma rota alternativa em que o
pagamento ocorrerá na direção oposta (zerando o saldo) e/ou
criando um caminho alternativo inteiramente para o pagamento. Isso criará
um valor temporal do dinheiro para divulgar entradas de hashes na Lightning
Network. Os participantes podem se especializar em alta conectividade
entre nós e oferecer-se para descarregar hashlocks de outros nós em troca
de uma taxa. Esses participantes concordarão com pagamentos que fazem net
out para zero (mais taxas), mas estão emprestando bitcoins por um período
de tempo definido. Provavelmente, essas entidades com baixa demanda por
recursos de canal serão usuários finais já conectados a múltiplos nós bem
conectados. Quando um usuário final se conecta a um nó, o nó pode pedir ao
cliente para bloquear seus fundos por vários dias em outro canal que o
cliente tenha estabelecido em troca de uma taxa. Isso pode ser obtido
fazendo com que as novas transações exijam um novo hash(Y) a partir da
entrada Y, além do hash existente, que pode ser gerado por qualquer
participante, mas deve divulgar Y apenas após estabelecer um círculo
completo. O novo participante tem a mesma responsabilidade e os mesmos
timelocks do antigo participante substituído. Também é possível que um
novo participante substitua múltiplos saltos.

\begin{figure}[H]
	\makebox[\linewidth]{
		\includegraphics[width=1\linewidth]{figures/network-fig4.pdf}
	}
	\caption{Erin está conectada tanto a Bob quanto a Dave. Se Bob deseja
		liberar seu canal com Carol, já que esse canal está ativo e
		muito lucrativo, Bob pode descarregar o pagamento para Dave
		via Erin. Como Erin tem bitcoin extra disponível, ela poderá
		coletar uma taxa por descarregar o canal entre Bob e Carol,
		bem como entre Carol e Dave. Os canais entre Bob e Carol,
		assim como entre Carol e Dave, são desfeitos e não têm mais
		o HTLC, nem ocorreu pagamento naquele caminho. O pagamento
		ocorrerá no caminho envolvendo Erin. Isso é alcançado
		criando um novo pagamento de Dave para Carol para Bob
		contingente à construção de um HTLC por Erin. Os pagamentos
		em linhas tracejadas (vermelho) são zerados e liquidados
		via um novo Contrato de Compromisso.
	}
\end{figure}

\subsection{Roteamento de Pagamentos}

É teoricamente possível construir um mapa de rotas implicitamente a partir
da observação de multisigs 2-de-2 na blockchain para construir uma tabela
de roteamento. Note, contudo, que isso não é viável com saídas de
transação pay-to-script-hash, que podem ser resolvidas fora de banda do
protocolo Bitcoin via um serviço de roteamento terceirizado. A construção
de uma tabela de roteamento se tornará necessária para grandes operadores
(por exemplo, BGP, Cjdns). Eventualmente, com otimizações, a rede se
parecerá muito com a rede bancária correspondente, ou com os ISPs Tier-1.
De forma similar à maneira como pacotes ainda chegam ao destino em uma
conexão de rede doméstica, nem todos os participantes precisam ter uma
tabela de roteamento completa. As rotas centrais Tier-1 podem estar
online o tempo todo \textemdash enquanto nós nas bordas, como usuários
comuns, estariam conectados de forma intermitente.

A descoberta de nós pode ocorrer ao longo das bordas pré-selecionando e
oferecendo rotas parciais para nós bem conhecidos.

\subsection{Taxas}

As taxas da Lightning Network, que diferem das taxas da blockchain, são
pagas diretamente entre participantes dentro do canal. As taxas pagam pelo
valor temporal do dinheiro consumindo o canal por um período máximo
determinado, e pelo risco de contraparte de não comunicação.

O risco de contraparte para taxas existe apenas com a contraparte direta
do canal. Se um nó a dois saltos de distância decide se desconectar e sua
transação é transmitida na blockchain, as contrapartes diretas não devem
transmitir na blockchain, mas continuar a atualizar via novação com uma
nova Transação de Compromisso. Veja a entrada Timelocks Decrescentes na
seção HTLC para mais informações sobre risco de contraparte.

O valor temporal das taxas paga pelo consumo de tempo (por exemplo, 3
dias) e é conceitualmente equivalente a uma taxa de leasing de ouro sem
risco custodial; é o valor temporal por consumir o acesso ao dinheiro por
uma duração muito curta. Como certos caminhos podem se tornar muito
lucrativos em uma direção, é possível que as taxas sejam negativas para
incentivar que o canal esteja disponível para esses caminhos lucrativos.

\section{Riscos}

Os principais riscos relacionam-se à expiração dos timelocks. Adicionalmente,
para que nós centrais e possivelmente alguns comerciantes possam rotear
fundos, as chaves precisam ser mantidas online para menor latência. No
entanto, usuários finais e nós podem manter suas chaves privadas isoladas em
carteira fria.

\subsection{Timelocks Inadequados}

Os participantes devem escolher timelocks com tempo suficiente. Se for dado
tempo insuficiente, é possível que transações com timelock que se acreditava
serem inválidas se tornem válidas, permitindo o roubo de moedas pela
contraparte. Existe um trade-off entre timelocks mais longos e o valor
temporal do dinheiro. Ao escrever software de carteira e de aplicações para
a Lightning Network, é necessário garantir que tempo suficiente seja dado e
que os usuários possam ter suas transações inseridas na blockchain ao
interagir com contrapartes não cooperativas ou maliciosas.

\subsection{Spam de Expiração Forçada}

A expiração forçada de muitas transações pode ser o maior risco sistêmico ao
usar a Lightning Network. Se um participante malicioso cria muitos canais e
força que todos expirem ao mesmo tempo, isso pode sobrecarregar a capacidade
de dados dos blocos, forçando expiração e transmissão à blockchain. O
resultado seria spam em massa na rede Bitcoin. O spam pode atrasar
transações ao ponto em que outras transações com locktime se tornem
válidas.

Isso pode ser mitigado permitindo uma substituição de transação em todas as
transações pendentes. Anti-spam pode ser usado permitindo apenas uma
substituição de transação de número de sequência maior pelo inverso de um
número par ou ímpar. Por exemplo, se um número de sequência ímpar foi
transmitido, permita uma substituição para um número par maior apenas uma
vez. As transações usariam o número de sequência de maneira ordenada para
substituir outras transações. Isso mitiga o risco assumindo mineradores
honestos. Este ataque é de risco extremamente alto, pois a transmissão
incorreta de Transações de Compromisso acarreta penalidade total de todos
os fundos no canal.

Adicionalmente, alguém pode tentar roubar transações HTLC forçando uma
transação de timeout a passar quando não deveria. Isso pode ser facilmente
mitigado fazendo com que cada transferência dentro do canal seja menor que
o total de taxas de transação utilizadas. Como as transações são
extremamente baratas e não atingem a blockchain com contrapartes de canal
cooperativas, grandes transferências de valor podem ser divididas em muitas
transferências pequenas. Essa tentativa só pode funcionar se os blocos
estiverem completamente cheios por um longo tempo. Embora seja possível
mitigá-la usando uma duração maior de timeout do HTLC, tamanhos de bloco
variáveis podem se tornar comuns, o que pode exigir mitigações.

Se esse tipo de transação se tornar a forma dominante de transações
incluídas na blockchain, pode se tornar necessário aumentar o tamanho do
bloco e executar uma estrutura de tamanho de bloco variável e flags de
timestop, conforme descrito na seção abaixo. Isso pode criar penalidades e
desincentivos suficientes para tornar a operação altamente não lucrativa e
malsucedida para os atacantes, já que os atacantes perdem todos os seus
fundos ao transmitir a transação errada, ao ponto em que isso jamais
ocorrerá.

\subsection{Roubo de Moedas via Quebra Criptográfica}

Como as partes precisam estar online e usar chaves privadas para assinar,
existe a possibilidade de que, se o computador onde as chaves privadas são
armazenadas for comprometido, as moedas serão roubadas pelo atacante.
Embora possa haver métodos para mitigar a ameaça para o remetente e o
destinatário, os nós intermediários devem estar online e provavelmente
processarão a transação automaticamente. Por essa razão, os nós
intermediários estarão em risco e não devem manter quantias substanciais de
dinheiro nessa ``carteira quente''. Nós intermediários com melhor segurança
provavelmente conseguirão superar outros no longo prazo e ser capazes de
conduzir maior volume de transações devido a taxas mais baixas.
Historicamente, um dos maiores componentes de taxas e juros no sistema
financeiro decorre de várias formas de risco de contraparte \textemdash no
Bitcoin é possível que o maior componente nas taxas seja derivado de prêmios
de risco de segurança.

Uma Transação de Financiamento pode ter múltiplas saídas com múltiplas
Transações de Compromisso, com a chave da Transação de Financiamento e
algumas chaves das Transações de Compromisso armazenadas offline. É
possível criar um equivalente a uma ``Conta Corrente'' e uma ``Conta
Poupança'' movendo fundos entre saídas de uma Transação de Financiamento,
com a ``Conta Poupança'' armazenada offline e exigindo assinaturas
adicionais de serviços de segurança.

\subsection{Perda de Dados}

Quando uma parte perde dados, é possível que a contraparte roube fundos.
Isso pode ser mitigado tendo um serviço de armazenamento de dados de
terceiros, no qual os dados criptografados são enviados a esse serviço
terceirizado que a parte não consegue descriptografar. Adicionalmente,
deve-se escolher contrapartes de canal que sejam responsáveis e dispostas a
fornecer o estado atual, com alguns testes periódicos de honestidade.

\subsection{Esquecer de Transmitir a Transação no Tempo Certo}

Se alguém não transmitir uma transação no momento correto, a contraparte
pode roubar fundos. Isso pode ser mitigado tendo um terceiro designado
para enviar fundos. Uma taxa de saída pode ser adicionada para criar um
incentivo para que esse terceiro monitore a rede. Além disso, isso também
pode ser mitigado implementando OP\_CHECKSEQUENCEVERIFY.

\subsection{Incapacidade de Realizar Soft-Forks Necessários}

Mudanças são necessárias no Bitcoin, como o soft-fork de maleabilidade.
Adicionalmente, se esse sistema se tornar popular, será necessário para o
sistema transacionar com segurança com muitos usuários, e algum tipo de
estrutura como um timestop por altura de bloco será desejável. Esse
sistema pressupõe tais mudanças para permitir que a Lightning Network
exista totalmente, bem como soft-forks que assegurem que a segurança seja
robusta contra atacantes. Embora o sistema possa continuar a operar
apenas com alguns soft-forks de timelock e maleabilidade, haverá
soft-forks necessários relativos a riscos sistêmicos. Sem a devida visão
da comunidade, a incapacidade de estabelecer um timestop ou função
similar permitirá que ataques sistêmicos ocorram e podem não ser
reconhecidos como imperativos até que um ataque realmente aconteça.

\subsection{Ataques de Mineradores em Conluio}

Mineradores podem optar por recusar incluir transações específicas (por
exemplo, transações de Remediação de Quebra) para auxiliar no roubo de
moedas por timeout. Um atacante pode pagar a todos os mineradores para
recusar a inclusão de certas transações em seu mempool e em seus blocos. Os
mineradores podem identificar seus próprios blocos na tentativa de provar
seu comportamento ao atacante pagador.

Isso pode ser mitigado incentivando os mineradores a evitar identificar
seus próprios blocos. Além disso, deve-se esperar que esse tipo de
pagamento a mineradores seja uma atividade maliciosa e que o contrato seja
inexequível. Os mineradores podem então aceitar o pagamento e minerar
furtivamente um bloco sem identificá-lo ao atacante. Como o atacante está
pagando por isso, ele rapidamente ficará sem dinheiro ao perder a taxa
para o minerador, bem como perder todo o seu dinheiro no canal. Esse
ataque é improvável e razoavelmente pouco atraente, pois é muito difícil e
requer um alto grau de conluio com risco extremo.

O modelo de risco desse ataque ocorrendo é similar ao de mineradores
conluiando para fazer ataques de reorg: extremamente improvável com muitos
mineradores não coordenados.

\section{Aumentos do Tamanho do Bloco e Consenso}

Se presumirmos que existe uma rede de pagamentos descentralizada e que um
usuário fará 3 transações na blockchain por ano em média, o Bitcoin poderá
suportar mais de 35 milhões de usuários com blocos de 1MB em circunstâncias
ideais (assumindo 2000 transações/MB, ou 500 bytes/Tx). Isso é bastante
limitado, e um aumento do tamanho do bloco pode ser necessário para
suportar todas as pessoas no mundo usando Bitcoin. Um simples aumento do
tamanho do bloco seria um hard-fork, ou seja, todos os nós precisariam
atualizar suas carteiras se desejassem participar da rede com os blocos
maiores.

Embora possa parecer que esse sistema mitigará os aumentos de tamanho de
bloco no curto prazo, se ele alcançar escala global, exigirá um aumento de
tamanho de bloco no longo prazo. Criar uma ferramenta crível para ajudar a
evitar spam na blockchain projetado para incentivar a expiração das
transações se torna imperativo.

Para mitigar vulnerabilidades de spam por timelock, as regras de consenso
para não-mineradores e mineradores também podem diferir, se as regras de
consenso dos mineradores forem mais restritivas. Os não-mineradores podem
aceitar blocos maiores que 1MB, enquanto os mineradores podem ter limites
diferentes de soft-cap nos tamanhos de bloco. Se um tamanho de bloco
estiver acima desse limite, então é visto como um bloco inválido por outros
mineradores, mas não por não-mineradores. Os mineradores apenas construirão
a cadeia sobre blocos que sejam válidos de acordo com o soft-cap acordado.
Isso permite que os mineradores concordem em elevar o limite de tamanho de
bloco sem exigir hard-forks frequentes dos clientes, contanto que o valor
elevado pelos mineradores não ultrapasse o limite rígido dos clientes. Isso
mitiga o risco de expiração em massa de transações de uma vez. Todas as
transações que não forem resgatadas via Liquidação por Exercício (ES) podem
ter uma taxa muito alta atrelada, e os mineradores podem usar uma regra de
consenso pela qual essas transações são isentas do soft-cap, tornando muito
provável que as transações corretas entrem na blockchain.

Quando as transações são vistas como circuitos e contratos em vez de
pacotes de transação, os riscos de consenso podem ser medidos pela
quantidade de tempo disponível para cobrir o conjunto de UTXO controlado
por partes hostis. Na prática, o limite superior do tamanho do UTXO é
determinado pelas taxas de transação e pelo valor mínimo padrão de saída
de transação. Se os mineradores do Bitcoin tiverem um mempool
determinístico que prioriza transações respeitando uma ordem temporal
local ``fraca'' das transações, pode se tornar extremamente não lucrativo e
improvável que um ataque seja bem-sucedido. Qualquer ataque de spam
temporal de transação transmitindo a Transação de Compromisso incorreta é
de risco extremamente alto para o atacante, pois requer uma imensa
quantidade de bitcoin e todos os fundos comprometidos nessas transações
serão perdidos se o atacante falhar.

\section{Casos de Uso}

Além de ajudar o Bitcoin a escalar, há muitos usos para as transações na
Lightning Network:
\begin{itemize}
	\item Transações Instantâneas. Usando a Lightning, as transações Bitcoin
		passam a ser quase instantâneas com qualquer parte. É possível
		pagar por uma xícara de café com pagamento direto e
		irrevogável em milissegundos a segundos.

	\item Arbitragem em Exchanges. Atualmente existe um incentivo para
		manter fundos em exchanges para estar pronto para grandes
		movimentos de mercado devido aos tempos de confirmação de
		3-6 blocos. É possível que a exchange participe dessa rede e
		que os clientes movam seus fundos para dentro e para fora da
		exchange para ordens quase instantaneamente. Se a exchange
		não tiver grande profundidade de mercado e se comprometer a
		permitir apenas ordens limite próximas ao topo do livro de
		ordens, então o risco de roubo de moedas se torna muito
		menor. A exchange, na prática, não teria mais necessidade de
		uma carteira fria. Isso pode reduzir
		substancialmente os roubos e a necessidade de custodiantes
		terceiros confiáveis.

	\item Micropagamentos. As taxas da blockchain do Bitcoin são altas
		demais para aceitar micropagamentos, especialmente com os
		menores valores. Com este sistema, micropagamentos quase
		instantâneos usando Bitcoin, sem um custodiante terceiro,
		seriam possíveis. Isso permitiria, por exemplo, pagar por
		megabyte por serviço de internet ou por artigo para ler um
		jornal.

	\item Contratos Financeiros Inteligentes e Escrow. Contratos financeiros
		são especialmente sensíveis ao tempo e têm maiores demandas
		por computação na blockchain. Ao mover a esmagadora maioria
		das transações sem confiança para fora da cadeia, é possível
		ter termos contratuais transacionais altamente complexos sem
		jamais atingir a blockchain.

	\item Pagamentos Cross-Chain. Contanto que haja funções de hash
		similares entre cadeias, é possível que as transações sejam
		roteadas em múltiplas cadeias com regras de consenso
		diferentes. O remetente não precisa confiar nem conhecer as
		outras cadeias \textemdash nem mesmo a cadeia de destino. De
		forma similar, o destinatário não precisa saber nada sobre a
		cadeia do remetente ou qualquer outra cadeia. Tudo o que
		interessa ao destinatário é um pagamento condicionado ao
		conhecimento de um segredo na sua cadeia. O pagamento pode
		ser roteado por participantes em ambas as cadeias no salto.
		Por exemplo, Alice está no Bitcoin, Bob está tanto no
		Bitcoin quanto na X-Coin, e Carol está em uma hipotética
		X-Coin; Alice pode pagar a Carol sem entender as regras de
		consenso da X-Coin.

\end{itemize}

\section{Conclusão}

Criar uma rede de canais de micropagamento permite a escalabilidade do
Bitcoin, micropagamentos até o nível do satoshi e transações quase
instantâneas. Esses canais representam transações Bitcoin reais, usando os
opcodes de scripting do Bitcoin para permitir a transferência de fundos
sem risco de roubo pela contraparte, especialmente com mitigações de
risco de mineradores de longo prazo.

Se todas as transações usando Bitcoin estivessem na blockchain, para
permitir que 7 bilhões de pessoas fizessem duas transações por dia, seriam
necessários, na melhor das hipóteses, blocos de 24GB a cada dez minutos
(assumindo 250 bytes por transação e 144 blocos por dia). Conduzir todas
as transações de pagamento globais na blockchain hoje implica que os
mineradores precisariam fazer uma quantidade incrível de computação,
limitando severamente a escalabilidade do Bitcoin e os nós completos a
alguns processadores centralizados.

Se todas as transações usando Bitcoin fossem conduzidas dentro de uma rede
de canais de micropagamento, para permitir que 7 bilhões de pessoas
fizessem dois canais por ano com transações ilimitadas dentro do canal,
seriam necessários blocos de 133 MB (assumindo 500 bytes por transação e
52.560 blocos por ano). Computadores desktop da geração atual poderão
executar um nó completo com blocos antigos podados em 2TB de armazenamento.

Com uma rede de canais de micropagamento instantaneamente confirmados
cujos pagamentos são onerados por timelocks e saídas com hashlocks, o
Bitcoin pode escalar para bilhões de usuários sem risco custodial ou
centralização da blockchain quando as transações são conduzidas com
segurança off-chain usando o scripting do Bitcoin, com a aplicação contra
a não cooperação feita pela transmissão de transações multiassinaturas
assinadas na blockchain.

\section{Agradecimentos}

Os canais de micropagamento foram desenvolvidos por muitas partes, e
foram discutidos no bitcointalk, na lista de e-mails do Bitcoin e no IRC.
A quantidade de contribuidores para essa ideia é imensa e muita reflexão
foi dedicada a essa capacidade. Foi feito um esforço para citar e
encontrar ideias similares, mas isso está absolutamente longe de estar
completo. Em particular, há muitas similaridades com uma proposta de
Alex Akselrod ao usar hashlocking como método de onerar um canal de
pagamento em estrela (hub-and-spoke).

Agradecimentos a Peter Todd por corrigir um erro significativo no script
do HTLC, bem como por otimizar o tamanho dos opcodes.

Agradecimentos a Elizabeth Stark pela revisão e correções.

Agradecimentos a Rusty Russell pela revisão deste documento e sugestões
para tornar o conceito mais digerível, bem como pelo trabalho em uma
construção que pode fornecer uma solução paliativa antes de uma correção
de maleabilidade de longo prazo (a ser descrita em uma versão futura).

\begin{appendices}
\section{Resolvendo a Maleabilidade}

Para criar esses contratos no Bitcoin sem um serviço de terceiros
confiável, o Bitcoin deve corrigir o problema de maleabilidade de
transações. Se as transações puderem ser mutadas, então as assinaturas
podem ser invalidadas, invalidando assim transações de reembolso e títulos
de compromisso. Isso cria uma oportunidade para atores hostis usarem-na
como uma oportunidade para uma tática de negociação para roubar moedas,
em essência, um cenário de refém.

Para mitigar a maleabilidade, é necessário fazer uma mudança via soft-fork
no Bitcoin. Clientes mais antigos ainda funcionarão, mas os mineradores
precisariam atualizar. O Bitcoin já teve diversos soft-forks no passado,
incluindo o pay-to-script-hash (P2SH).

Para mitigar a maleabilidade, é necessário alterar quais conteúdos são
assinados pelos participantes. Isso é alcançado criando novos tipos de
sighash. Para acomodar esse novo comportamento, um novo tipo de P2SH ou
um novo OP\_CHECKSIG é necessário para torná-lo um soft-fork em vez de um
hard-fork.

Se um novo P2SH fosse definido, ele usaria um script de saída diferente,
como:
\\\\
\noindent\texttt{OP\_DUP OP\_HASH160 \textless hash de 20 bytes\textgreater
~OP\_EQUALVERIFY}
\\\\

Como isso sempre resolverá para verdadeiro contanto que um redeemScript
válido seja fornecido, todos os clientes existentes retornarão verdadeiro.
Isso permite que o sistema de scripting construa novas regras, incluindo
novas regras de validação de assinatura. Pelo menos um novo sighash
precisaria existir.

SIGHASH\_NOINPUT não assinaria nenhum ID de transação de entrada nem
assinaria o índice. Ao usar SIGHASH\_NOINPUT, pode-se garantir que a
contraparte não invalide árvores inteiras de transações encadeadas de
potenciais estados de contrato previamente acordados, via mutação do ID
da transação. Com os novos flags de sighash, é possível gastar a partir de
uma transação pai mesmo que o ID da transação tenha mudado, contanto que o
script avalie como verdadeiro (ou seja, uma assinatura válida).

SIGHASH\_NOINPUT implica em risco significativo com a reutilização de
endereços, pois pode funcionar com qualquer transação em que o sigScript
retorne como válido, de modo que várias transações com as mesmas saídas são
resgatáveis (contanto que os valores de saída sejam menores).

Além disso, e tão importante quanto, SIGHASH\_NOINPUT permite que os
participantes assinem gastos de transações sem conhecer as assinaturas da
transação sendo gasta. Resolvendo a maleabilidade da maneira acima, duas
partes podem construir contratos e gastar transações sem que qualquer das
partes tenha a capacidade de transmitir essa transação original na
blockchain até que ambas concordem. Com o novo tipo de sighash, os
participantes podem construir possíveis estados de contrato e possíveis
condições de pagamento e concordar com todos os termos, antes que o
contrato possa ser pago, transmitido e executado, sem a necessidade de um
terceiro confiável.

Sem SIGHASH\_NOINPUT, não é possível construir saídas antes que a
transação possa ser financiada. É como se não se pudesse fazer acordos sem
comprometer fundos sem saber a que se está comprometendo.
SIGHASH\_NOINPUT permite construir o resgate de transações que ainda não
existem. Em outras palavras, é possível formar acordos antes de financiar
a transação, se a saída for uma transação multiassinatura 2-de-2.

Para usar SIGHASH\_NOINPUT, constrói-se uma Transação de Financiamento e
não a assina ainda. Essa Transação de Financiamento não precisa usar
SIGHASH\_NOINPUT se estiver gastando a partir de uma transação que já
entrou na blockchain. Para gastar a partir de uma Transação de
Financiamento com uma saída multiassinatura 2-de-2 que ainda não foi
assinada e transmitida, no entanto, é necessário usar SIGHASH\_NOINPUT.

Uma solução paliativa adicional usando OP\_CHECKSEQUENCEVERIFY ou um uso
menos ideal de OP\_CHECKLOCKTIMEVERIFY será descrita em um artigo futuro
de Rusty Russell. Uma versão atualizada deste artigo também incluirá essas
construções.

\end{appendices}

%bibliografia
\bibliographystyle{unsrt}
\bibliography{paper}

\end{document}

%TODO Remover "Fidelity Bond" ("caução de fidelidade") e substituir por um nome
%que se refira a ele como um contrato, em vez disso.

